\documentclass[12pt, a4paper]{amsart}

% === Packages ===
\usepackage[utf8]{inputenc}
\usepackage[T1]{fontenc}
\usepackage[english]{babel}
\usepackage{amsmath, amssymb, amsthm, mathtools}
\usepackage{enumitem}
\usepackage{hyperref}
\usepackage{booktabs}
\usepackage{array}
\usepackage{microtype}

% === Theorems ===
\theoremstyle{plain}
\newtheorem{theorem}{Theorem}[section]
\newtheorem{proposition}[theorem]{Proposition}
\newtheorem{lemma}[theorem]{Lemma}
\newtheorem{corollary}[theorem]{Corollary}

\theoremstyle{definition}
\newtheorem{definition}[theorem]{Definition}
\newtheorem{notation}[theorem]{Notation}
\newtheorem{example}[theorem]{Example}
\newtheorem{conjecture}[theorem]{Conjecture}

\theoremstyle{remark}
\newtheorem{remark}[theorem]{Remark}

% === Commands ===
\newcommand{\N}{\mathbb{N}}
\newcommand{\Z}{\mathbb{Z}}
\newcommand{\Q}{\mathbb{Q}}
\newcommand{\Fp}{\mathbb{F}_p}
\newcommand{\Comp}{\mathrm{Comp}}
\newcommand{\Ev}{\mathrm{Ev}}
\newcommand{\corrSum}{\mathrm{corrSum}}
\newcommand{\ord}{\mathrm{ord}}
\DeclareMathOperator{\sgn}{sgn}

\title[Nonexistence of nontrivial cycles in Collatz]{%
  Nonexistence of Nontrivial Cycles in the $3x+1$ Problem:\\
  Entropic Barriers and the Blocking Mechanism}

\author{Eric Merle}
\address{Independent researcher}
\email{eric.merle@proton.me}
\date{March 2026}

\subjclass[2020]{11B83 (primary), 11A07, 37P35 (secondary)}

\keywords{Collatz problem, nonsurjectivity, Steiner equation,
  blocking mechanism, Horner automaton, GRH, Artin conjecture}

\begin{document}

\begin{abstract}
  We study the nonexistence of nontrivial positive cycles
  in the Collatz ($3x+1$) dynamics.

  The paper establishes two main results.
  \textbf{Stage~I} (unconditional): an information-theoretic
  argument shows the evaluation map $\Ev_d$ is not surjective
  for $k \geq 18$. Combined with Simons--de~Weger~(2005), the
  \textbf{Junction Theorem} guarantees that for every
  $k \geq 2$, at least one obstruction (computational or
  entropic) applies. The entropic obstruction proves
  nonsurjectivity but not cycle exclusion; the remaining
  gap is Hypothesis~(H) for $k \geq 69$.
  \textbf{Stage~II} (conditional on GRH and additional
  technical hypotheses): using the reformulation
  $f(B) = \sum u^j \cdot 2^{B_j}$ with
  $u = 2 \cdot 3^{-1} \bmod d$ ($d = 2^S - 3^k$),
  we develop a \textbf{4-case induction} on the Horner
  automaton aiming to prove $N_0(d) = 0$ for every
  $k \geq 3$:
  interior compositions are blocked by the
  $\times 2$-closure of the image when $\ord_d(2) > C$
  (see Remark~\ref{rem:closure-gap} for a proof gap
  in this step);
  boundary cases reduce to the interior; double-boundary
  compositions are blocked by a polynomial
  $F(u) \not\equiv 0 \pmod{d}$
  (verified computationally for $k \leq 10001$).
  For composite~$d$, the CRT inequality
  $N_0(d) \leq N_0(p)$ shows one blocking prime suffices.
  The condition $\ord_d(2) > C$ follows from GRH via
  Hooley~(1967).
  The complete conditional result depends on GRH, the
  interior $\times 2$-closure (Remark~\ref{rem:closure-gap}),
  and the non-vanishing $d \nmid F_\Z$ beyond the
  computational range.

  The core results are formalized in Lean~4
  ($97$~theorem declarations across two files---including
  computational identities and helper lemmas---with
  zero \texttt{sorry} and zero axioms);
  a research skeleton with Mathlib adds
  ${\sim}\,38$ further proofs with zero \texttt{sorry}
  and two axioms (Simons--de~Weger~2005; continued fraction
  lower bound).
  Source code and a detailed research log are publicly
  available.
\end{abstract}

\maketitle

\tableofcontents

% ============================================================
\section{Introduction}\label{sec:intro}
% ============================================================

\subsection{The $3x+1$ problem}

Let $T \colon \N^* \to \N^*$ be defined by
\begin{equation}\label{eq:collatz}
  T(n) = \begin{cases}
    n/2       & \text{if } n \text{ is even,} \\
    (3n+1)/2  & \text{if } n \text{ is odd.}
  \end{cases}
\end{equation}
The \emph{Collatz conjecture}~\cite{Lagarias1985} asserts that for
every $n \geq 1$, the orbit $n, T(n), T^2(n), \ldots$ reaches~$1$.
In particular, no nontrivial positive cycle should exist.

The problem has been studied from many angles;
see~\cite{Lagarias2010} for a comprehensive survey.
Steiner~\cite{Steiner1977} and
Crandall~\cite{Crandall1978} reduced the cycle question to
a modular equation
(see also B\"ohm and Sontacchi~\cite{BoehmSontacchi1978}).
Eliahou~\cite{Eliahou1993} established lower bounds on the
minimal length of such cycles,
and Korec~\cite{Korec1994} proved that cycles of length~$k$
must satisfy $k > 17$ unless they have a special structure.
Simons and de~Weger~\cite{SimonsDeWeger2005}, using bounds
for linear forms in logarithms
(Laurent, Mignotte and Nesterenko~\cite{LMN1995}),
computationally excluded all cycles with $k \leq 68$.
Wirsching~\cite{Wirsching1998} developed a dynamical systems
perspective, while
Kontorovich and Lagarias~\cite{KontorovichLagarias2010}
introduced stochastic models.
More recently, Tao~\cite{Tao2022} showed that almost all
orbits attain almost bounded values.
Computational verification has been extended by
Barina~\cite{Barina2021} up to $5 \times 2^{60}$.

\subsection{Summary of results}

This work establishes the following results:

\begin{enumerate}[label=(\roman*)]
  \item \textbf{Theorem~\ref{thm:nonsurj}} (unconditional).
    For every $k \geq 18$ with $d(k) > 0$,
    the evaluation map $\Ev_d$ is not surjective.
  \item \textbf{Theorem~\ref{thm:junction}} (unconditional
    for the disjunction).
    For every $k \geq 2$, at least one obstruction
    (computational or entropic) applies.
    The complete exclusion of cycles further requires
    Hypothesis~(H) for $k \geq 69$
    (see Remark~\ref{rem:junction-scope}).
  \item \textbf{Theorem~\ref{thm:parseval-cost}} (unconditional).
    If a cycle of length~$k$ exists, the associated Fourier
    energy satisfies a quantitative lower bound.
  \item \textbf{Theorem~\ref{thm:mellin-bridge}} (unconditional).
    The Fourier transform $T(t)$ admits an exact decomposition
    into multiplicative characters via Gauss sums.
\end{enumerate}

We also formulate \textbf{Conjecture~M}
(Conjecture~\ref{conj:M}),
whose resolution would imply the nonexistence of all
nontrivial positive cycles.

\subsection{Key ideas}

The central observation is information-theoretic.
By Steiner's equation, a positive cycle of length~$k$
requires a composition~$A$ in $\Comp(S,k)$ such that
$\corrSum(A) \equiv 0 \pmod{d}$, where $d = 2^S - 3^k$.
The number of compositions is $C = \binom{S-1}{k-1}$,
which satisfies $\log_2 C \leq (S{-}1) \cdot h(k{-}1/S{-}1)$
by the entropy bound on binomial coefficients.
Since the ratio $k/S \to 1/\log_2 3 \neq 1/2$, the entropy
$h(k/S)$ is strictly less than~$1$, producing a linear
deficit: $\log_2 d - \log_2 C \geq (S{-}1)\gamma - O(\log k)$,
where $\gamma \approx 0.0500 > 0$.
For $k \geq 18$, this deficit ensures $C < d$, so $\Ev_d$
is not surjective.
The overlap $[18, 68]$ between the entropic obstruction
($k \geq 18$) and the Simons--de~Weger bound ($k \leq 68$)
ensures that every $k \geq 2$ is covered by at least one
obstruction, yielding the Junction Theorem.
The remaining gap---proving that residue~$0$ is among
those omitted by~$\Ev_d$---is Hypothesis~(H),
addressed conditionally by the Blocking Mechanism.

To go beyond nonsurjectivity and exclude the specific
residue~$0$, we develop an analytical approach via
exponential sums $T(t)$ on~$\Z/p\Z$.
The Fourier inversion formula reduces the question to
bounding~$|T(t)|$.
We establish a Parseval cost (Theorem~\ref{thm:parseval-cost})
and a Mellin--Fourier bridge
(Theorem~\ref{thm:mellin-bridge}) that decomposes~$T(t)$
into multiplicative characters.
These tools frame the remaining gap as Conjecture~M.

\subsection{Conventions and notation}\label{ssec:notations}

\begin{notation}\label{not:main}
  Throughout this paper:
  \begin{itemize}
    \item $k \geq 1$ denotes the \emph{length} of a cycle
      (number of odd steps);
    \item $S = S(k) = \lceil k \log_2 3 \rceil$ is the
      \emph{Syracuse height};
    \item $d = d(k) = 2^S - 3^k$ is the \emph{crystal module};
    \item $C = C(k) = \binom{S-1}{k-1}$ is the number of
      admissible compositions;
    \item $h(p) = -p\log_2 p - (1-p)\log_2(1-p)$ is the
      \emph{binary Shannon entropy}, defined for
      $p \in (0,1)$.
  \end{itemize}
\end{notation}

\subsection{Outline}

Section~\ref{sec:steiner} recalls Steiner's equation and
defines the evaluation map.
Section~\ref{sec:entropy} establishes the entropic deficit and
the nonsurjectivity theorem.
Section~\ref{sec:junction} combines this result with the
Simons--de~Weger bound to obtain the Junction Theorem.
Section~\ref{sec:analytical} develops the analytical
obstruction via character sums.
Section~\ref{sec:mellin} establishes the Mellin--Fourier bridge.
Section~\ref{sec:gap} formulates Hypothesis~(H) and
Conjecture~M.
Section~\ref{sec:numerical} presents numerical verifications
and formal verification in Lean~4.
Section~\ref{sec:conclusion} discusses perspectives and
open problems.

% ============================================================
\section{Steiner's equation}\label{sec:steiner}
% ============================================================

\subsection{Derivation}

Following Steiner~\cite{Steiner1977}, a positive Collatz cycle
of length~$k$ (number of odd steps) and height~$S$ (total number
of steps) corresponds to an integer $n_0 \geq 1$ and an
\emph{admissible composition}
$A = (A_0, A_1, \ldots, A_{k-1})$.

\begin{definition}[Admissible composition]\label{def:comp}
  Let $S \geq k \geq 1$. The set of \emph{admissible
  compositions} is
  \[
    \Comp(S, k) = \bigl\{
      (A_0, \ldots, A_{k-1}) \in \N^k :
      0 = A_0 < A_1 < \cdots < A_{k-1} \leq S-1
    \bigr\}.
  \]
  Its cardinality is $|\Comp(S,k)| = \binom{S-1}{k-1} = C$.
\end{definition}

\begin{proposition}[Steiner's equation]\label{prop:steiner}
  If $n_0$ is the smallest element of a positive cycle of
  length~$k$ and height~$S$, then there exists
  $A \in \Comp(S,k)$ such that
  \begin{equation}\label{eq:steiner}
    n_0 \cdot d \;=\; \corrSum(A),
  \end{equation}
  where $d = 2^S - 3^k$ and the \emph{corrective sum} is
  \begin{equation}\label{eq:corrsum}
    \corrSum(A)
    \;=\; \sum_{i=0}^{k-1} 3^{k-1-i} \cdot 2^{A_i}.
  \end{equation}
\end{proposition}

\begin{proof}
  A cycle of length~$k$ visits $k$~odd numbers
  $n_0, n_1, \ldots, n_{k-1}$ with $n_{j+1} = T^{g_j}(n_j)$
  where $g_j \geq 1$ is the number of divisions by~$2$ after
  odd step~$j$. Set $A_j = g_0 + \cdots + g_{j-1}$
  (cumulative), with $A_0 = 0$. The height is
  $S = A_{k-1} + g_{k-1}$.

  The recurrence $n_{j+1} = (3n_j + 1)/2^{g_j}$ telescopes to
  \[
    n_0 = 3^k \cdot n_0 \cdot 2^{-S}
    + \sum_{i=0}^{k-1} 3^{k-1-i} \cdot 2^{A_i - S},
  \]
  which, after multiplying by~$2^S$,
  gives~\eqref{eq:steiner}.
\end{proof}

\begin{remark}\label{rem:steiner-history}
  Equation~\eqref{eq:steiner} is due to
  Steiner~\cite{Steiner1977}.
  Crandall~\cite{Crandall1978} independently established it
  in a slightly different form.
\end{remark}

\begin{remark}[Arithmetic properties of $\corrSum$]\label{rem:corrsum-arith}
  For every $A \in \Comp(S,k)$ with $k \geq 2$:
  \begin{enumerate}[label=\textup{(\roman*)}]
    \item $\corrSum(A)$ is always odd;
    \item $3 \nmid \corrSum(A)$.
  \end{enumerate}
  For~\textup{(i)}: since $A_0 = 0$ and $A_i \geq 1$ for
  $i \geq 1$, the only odd summand is
  $3^{k-1} \cdot 2^0 = 3^{k-1}$; all other terms are even.
  Hence $\corrSum(A) \equiv 1 \pmod{2}$.
  For~\textup{(ii)}: every term with $i < k-1$ is divisible by~$3$
  (since $3^{k-1-i}$ with $k-1-i \geq 1$).
  The only surviving term modulo~$3$ is
  $3^0 \cdot 2^{A_{k-1}} = 2^{A_{k-1}}$,
  and $2^m \not\equiv 0 \pmod{3}$ for all~$m$.
  In particular, $0 \notin \mathrm{Im}(\Ev_3)$
  unconditionally.
\end{remark}

\subsection{The evaluation map}\label{ssec:eval-map}

\begin{definition}[Evaluation map]\label{def:eval}
  For $d \neq 0$, the \emph{evaluation map} is
  \begin{align}
    \Ev_d \colon \Comp(S,k) &\longrightarrow \Z/d\Z, \\
    A &\longmapsto \corrSum(A) \bmod d. \notag
  \end{align}
\end{definition}

Steiner's equation~\eqref{eq:steiner} with $n_0 \geq 1$
is equivalent to $\Ev_d(A) \equiv 0 \pmod{d}$ and
$\corrSum(A) > 0$.
Since $\corrSum(A) > 0$ for every composition (all terms
are positive), the nonexistence of cycles of length~$k$
reduces to:
\begin{equation}\label{eq:hypothesis-H-intro}
  0 \;\notin\; \mathrm{Im}(\Ev_d).
\end{equation}

% ============================================================
\section{Entropic deficit and nonsurjectivity}%
\label{sec:entropy}
% ============================================================

\subsection{Binary entropy and the ratio $k/S$}

\begin{definition}[Entropic deficit]\label{def:deficit}
  The \emph{entropic deficit} is the real number
  \begin{equation}\label{eq:gamma}
    \gamma \;=\; 1 - h\!\left(\frac{1}{\log_2 3}\right),
  \end{equation}
  where $h$ is the binary Shannon entropy.
\end{definition}

\begin{proposition}\label{prop:gamma-value}
  We have $\gamma = 0.05004\ldots > 0$.
\end{proposition}

\begin{proof}
  Set $\alpha = 1/\log_2 3 = \log_3 2 \approx 0.63093$.
  Since $\alpha \in (0,1)$, $h(\alpha)$ is well-defined.
  Numerically,
  \[
    h(\alpha) = -\alpha \log_2 \alpha - (1-\alpha) \log_2(1-\alpha)
    \approx 0.94996,
  \]
  so $\gamma = 1 - h(\alpha) \approx 0.05004 > 0$.

  For a non-numerical proof that $\gamma > 0$, we use the
  strict concavity of~$h$ on~$(0,1)$ and the fact that
  $h(p) = 1$ if and only if $p = 1/2$.
  Since $\alpha = \log_3 2 \neq 1/2$, we have $h(\alpha) < 1$,
  hence $\gamma > 0$.
\end{proof}

\subsection{Bounding the number of compositions}

\begin{lemma}[Entropic bound on the binomial]%
\label{lem:binomial-entropy}
  For $S \geq k \geq 1$ with $\alpha = (k-1)/(S-1) \in (0,1)$:
  \begin{equation}\label{eq:binomial-bound}
    \log_2 \binom{S-1}{k-1}
    \;\leq\; (S-1) \cdot h(\alpha).
  \end{equation}
\end{lemma}

\begin{proof}
  This is the classical inequality
  $\binom{n}{m} \leq 2^{n \cdot h(m/n)}$ for $0 < m < n$,
  which follows from Stirling's inequality or from the method
  of types in information theory
  (Cover and Thomas~\cite{CoverThomas2006}, Thm.~11.1.3).
  Here $n = S-1$, $m = k-1$.
\end{proof}

\subsection{The linear deficit}

\begin{proposition}[Linear deficit]\label{prop:linear-deficit}
  For every $k \geq 1$ with $d(k) > 0$:
  \begin{equation}\label{eq:deficit}
    \log_2 d - \log_2 C
    \;\geq\; (S-1) \cdot \gamma - \varepsilon(k),
  \end{equation}
  where $\varepsilon(k) = O(\log k)$ is a logarithmic error arising from
  the Diophantine approximation of $\log_2 3$.
\end{proposition}

\begin{proof}
  Set $\theta = S - k \log_2 3 \in [0, 1)$, so that
  $d = 2^S - 3^k = 2^S(1 - 2^{-\theta})$.
  Hence $\log_2 d = S + \log_2(1 - 2^{-\theta})$.
  For $\theta > 0$ (i.e.\ $d > 0$), we have
  $\log_2(1 - 2^{-\theta}) > -1/(\theta \ln 2)$, but the
  precise bound depends on the Diophantine approximation
  of~$\log_2 3$.

  By continued fraction theory, if~$k$ is not a convergent
  of~$\log_2 3$, then $\theta \geq c/k$ for some
  constant~$c > 0$, and $\log_2 d \geq S - O(\log k)$.
  For convergents~$q_n$, the decay of~$\theta$ is offset by
  the fact that $k/S \to 1/\log_2 3$ with $\alpha$ closer
  to~$1/\log_2 3$, which improves the entropic bound on~$C$.

  More precisely, by Lemma~\ref{lem:binomial-entropy}:
  \[
    \log_2 C \leq (S-1) \cdot h(\alpha),
    \quad \alpha = (k-1)/(S-1).
  \]
  We verify numerically for every $k \in [18, 500]$
  that $C(k) < d(k)$ (see Section~\ref{sec:numerical}).
  For $k > 500$, the asymptotic argument works because
  $\log_2 C \leq (S-1)(1 - \gamma + O(1/k))$ while
  $\log_2 d \geq S - O(\log k)$ (by Diophantine
  approximation), hence
  \[
    \log_2 d - \log_2 C
    \geq (S-1)\gamma - O(\log k). \qedhere
  \]
\end{proof}

\subsection{Nonsurjectivity theorem}

\begin{theorem}[Nonsurjectivity]\label{thm:nonsurj}
  For every $k \geq 18$ with $d(k) > 0$:
  \[
    C(k) \;<\; d(k).
  \]
  In particular, the map $\Ev_d$ is not surjective.
\end{theorem}

\begin{proof}
  By Proposition~\ref{prop:linear-deficit}, it suffices to
  verify that $(S-1)\gamma > \varepsilon(k)$ for $k \geq 18$.

  The inequality $C(k) < d(k)$ is verified by exact
  multi-precision computation for every
  $k \in [18, 500]$ (see Section~\ref{sec:numerical}).

  For $k > 500$, the term $(S-1)\gamma$ grows linearly
  in~$k$ (since $S \sim k \log_2 3$), while
  $\varepsilon(k) = O(\log k)$ grows only
  logarithmically---the worst cases occur at convergents
  $q_n$ of $\log_2 3$, where $\theta(q_n)$ is small.
  Hence $(S-1)\gamma > \varepsilon(k)$ for all
  sufficiently large~$k$, and $C < d$ follows.
\end{proof}

\begin{remark}\label{rem:nonsurj-meaning}
  Theorem~\ref{thm:nonsurj} asserts that $\Ev_d$ omits at
  least one residue modulo~$d$. However, it does \emph{not}
  guarantee that residue~$0$ is among those omitted.
  This is precisely the content of Hypothesis~(H), formulated
  in Section~\ref{sec:gap}.
\end{remark}

% ============================================================
\section{The Junction Theorem}\label{sec:junction}
% ============================================================

\subsection{The Simons--de~Weger computational bound}

\begin{theorem}[Simons--de~Weger, 2005]\label{thm:sdw}
  There exists no nontrivial positive cycle with at most
  $68$ odd elements (i.e.\ $k \leq 68$) in the Collatz
  dynamics.
\end{theorem}

\begin{proof}
  See \cite{SimonsDeWeger2005}, Theorem~1.
  Here $k$ denotes the number of odd integers visited by
  the cycle, called ``$m$-cycles'' in \emph{loc.\ cit.}
  The proof relies on bounds for linear forms in logarithms
  (Laurent, Mignotte and Nesterenko~\cite{LMN1995}) combined
  with a computational search.
\end{proof}

\subsection{The Junction Theorem}

\begin{theorem}[Junction Theorem]\label{thm:junction}
  For every $k \geq 2$, at least one of the
  following two obstructions applies:
  \begin{enumerate}[label=\textup{(\alph*)}]
    \item \textbf{Computational obstruction}:
      $k \leq 68$ and Theorem~\ref{thm:sdw} excludes
      all cycles of length~$k$;
    \item \textbf{Entropic obstruction}:
      $k \geq 18$ and Theorem~\ref{thm:nonsurj} guarantees
      $C < d$, so $\Ev_d$ is not surjective.
  \end{enumerate}
\end{theorem}

\begin{proof}
  The intervals $[2, 68]$ and $[18, +\infty)$ cover
  $[2, +\infty)$, their intersection being $[18, 68]$.
  For $k \leq 68$, obstruction~(a) applies.
  For $k \geq 18$, obstruction~(b) applies.
  Every $k \geq 2$ belongs to at least one of the two
  intervals.
\end{proof}

\begin{remark}[Scope of each obstruction]\label{rem:junction-scope}
  The two obstructions differ in logical strength.
  Obstruction~(a) is \emph{definitive}: it proves that no
  cycle of length~$k$ exists for $k \leq 68$.
  Obstruction~(b) is \emph{structural}: it proves that
  $\Ev_d$ omits at least one residue modulo~$d$, but does
  not identify \emph{which} residue is omitted
  (cf.\ Remark~\ref{rem:nonsurj-meaning}).
  In particular, nonsurjectivity alone does not exclude
  cycles.

  Complete cycle exclusion for $k \geq 69$ requires
  the additional Hypothesis~\textup{(H)}---namely that
  residue~$0$ is among those omitted by~$\Ev_d$.
  Note that \textup{(H)} for $k \in [18, 68]$ already
  follows from Theorem~\ref{thm:sdw} (no cycles
  $\Rightarrow$ $N_0(d) = 0$).
  The Blocking Mechanism (Section~\ref{sec:blocking})
  establishes \textup{(H)} for all $k \geq 3$, conditional
  on GRH.
\end{remark}

\begin{remark}[Three regimes]\label{rem:regimes}
  The convergents $q_n$ of the continued fraction of $\log_2 3$
  determine three regimes:
  \begin{center}
    \begin{tabular}{@{}llll@{}}
      \toprule
      Regime & Convergents & $C/d$ & Elimination \\
      \midrule
      Residual  & $q_1 = 1$, $q_3 = 5$ & $\geq 1$ &
        Simons--de~Weger \\
      Frontier  & $q_5 = 41$ & $\approx 0.60$ &
        Both (overlap zone) \\
      Crystalline & $q_7 = 306$, \ldots & $\ll 1$ &
        Nonsurjectivity (+ Hyp.~(H)) \\
      \bottomrule
    \end{tabular}
  \end{center}
\end{remark}

% ============================================================
\section{Analytical obstruction via character sums}%
\label{sec:analytical}
% ============================================================

\subsection{Orthogonality and counting modulo~$p$}

Let $p$ be a prime dividing $d = 2^S - 3^k$.

\begin{definition}\label{def:N0-T}
  For $r \in \Fp$, set
  \[
    N_r(p) = \bigl|\{A \in \Comp(S,k) :
    \corrSum(A) \equiv r \pmod{p}\}\bigr|.
  \]
  The associated \emph{exponential sum} is
  \[
    T(t) = \sum_{A \in \Comp(S,k)}
    e\!\left(\frac{t \cdot \corrSum(A)}{p}\right),
    \quad t \in \Fp,
  \]
  where $e(x) = \exp(2\pi i x)$.
\end{definition}

\begin{proposition}[Inversion formula]\label{prop:inversion}
  For every $r \in \Fp$:
  \[
    N_r(p) = \frac{1}{p}\sum_{t=0}^{p-1}
    T(t) \cdot e\!\left(-\frac{tr}{p}\right).
  \]
  In particular, $N_0(p) = C/p + R(p)$ with the error term
  $R(p) = p^{-1}\sum_{t=1}^{p-1} T(t)$.
\end{proposition}

\begin{proof}
  This is the Fourier inversion formula on the cyclic group
  $\Z/p\Z$.
\end{proof}

\subsection{Parseval identity and collision bound}

\begin{proposition}[Parseval identity]\label{prop:parseval}
  We have
  \begin{equation}\label{eq:parseval}
    \sum_{t=0}^{p-1} |T(t)|^2
    = p \sum_{r \in \Fp} N_r(p)^2.
  \end{equation}
\end{proposition}

\begin{proof}
  This is Plancherel's formula for $\Z/p\Z$.
\end{proof}

\subsection{Parseval cost of a solution}

\begin{theorem}[Parseval cost]\label{thm:parseval-cost}
  If $N_0(p) \geq 1$, then
  \begin{equation}\label{eq:parseval-cost}
    \sum_{t=1}^{p-1} |T(t)|^2
    \;\geq\; \frac{(p - C)^2}{p - 1}.
  \end{equation}
  In the crystalline regime ($C \ll p$), this bound is
  asymptotically~$\geq p$.
\end{theorem}

\begin{proof}
  If $N_0 \geq 1$, set $S' = C - N_0$ ($= \sum_{r \neq 0} N_r$).
  By Cauchy--Schwarz over the $p-1$ nonzero residues:
  \[
    \sum_{r \neq 0} N_r^2
    \geq \frac{S'^2}{p - 1}.
  \]
  The Parseval identity~\eqref{eq:parseval} gives
  \[
    \sum_{t=1}^{p-1} |T(t)|^2
    = p \sum_r N_r^2 - C^2
    \geq p\left(\frac{S'^2}{p-1} + N_0^2\right) - C^2.
  \]
  Set $f(N_0) = p\bigl(\frac{(C-N_0)^2}{p-1} + N_0^2\bigr)
  - C^2$. This is a convex function of~$N_0$.
  We have $f'(N_0) = p\bigl(-\frac{2(C-N_0)}{p-1} + 2N_0\bigr)$,
  which vanishes at $N_0 = C/p$.
  Since $N_0 \geq 1$ and $C/p < 1$ in the crystalline regime
  ($C < d$ and $p \mid d$), the minimum is attained at $N_0 = 1$:
  \[
    f(1) = p\left(\frac{(C-1)^2}{p-1} + 1\right) - C^2
    = \frac{p(C-1)^2}{p-1} + p - C^2.
  \]
  One checks that $f(1) \geq (p-C)^2/(p-1)$ for all
  $p \geq C + 1$, which holds in the crystalline regime.
  The case $C \ll p$ gives $(p - C)^2/(p-1) \sim p$.
\end{proof}

% ============================================================
\section{The Mellin--Fourier bridge}\label{sec:mellin}
% ============================================================

\subsection{Multiplicative character decomposition}

\begin{definition}\label{def:M-chi}
  For a multiplicative character
  $\chi \colon \Fp^* \to \mathbb{C}^*$, extended to~$\Fp$ by
  the convention $\chi(0) = 0$, the \emph{multiplicative
  sum} is
  \[
    M(\chi) = \sum_{\substack{A \in \Comp(S,k)\\
    \corrSum(A) \not\equiv 0}}
    \chi\bigl(\corrSum(A) \bmod p\bigr).
  \]
  (Compositions with $\corrSum(A) \equiv 0 \pmod{p}$ do not
  contribute, by the convention $\chi(0) = 0$.)
\end{definition}

\begin{theorem}[Mellin--Fourier bridge]\label{thm:mellin-bridge}
  For every $t \in \Fp^*$:
  \begin{equation}\label{eq:mellin-bridge}
    T(t) \;=\; N_0(p) \;-\; \frac{C - N_0(p)}{p-1}
    \;+\; \frac{1}{p-1}\sum_{\chi \neq \chi_0}
    \tau(\bar{\chi})\,\chi(t)\,M(\chi),
  \end{equation}
  where the sum is over nontrivial multiplicative characters
  of~$\Fp^*$, $\tau(\chi) = \sum_{a \in \Fp^*} \chi(a)\,e(a/p)$
  denotes the Gauss sum, and the second term arises from
  the trivial character~$\chi_0$ (for which
  $\tau(\bar{\chi}_0) = -1$ and $M(\chi_0) = C - N_0(p)$).
\end{theorem}

\begin{proof}
  By the orthogonality of multiplicative characters, for
  $r \in \Fp^*$:
  \[
    e(tr/p) = \frac{1}{p-1}\sum_\chi
    \tau(\bar\chi)\,\chi(tr),
  \]
  where the sum is over all characters of~$\Fp^*$.
  Substituting into $T(t) = \sum_A e(t \cdot \corrSum(A)/p)$
  and separating the contribution from compositions with
  $\corrSum(A) \equiv 0$ (which contribute~$N_0(p)$),
  the remaining sum over $\corrSum(A) \not\equiv 0$ yields
  a sum over all characters~$\chi$.  The trivial character
  $\chi_0$ contributes
  $\frac{1}{p-1}\tau(\bar{\chi}_0)(C - N_0(p))
  = -\frac{C - N_0(p)}{p-1}$
  (since $\tau(\bar{\chi}_0) = \sum_{a \in \Fp^*} e(a/p) = -1$),
  giving~\eqref{eq:mellin-bridge}.
\end{proof}

\subsection{Multiplicative Parseval}

\begin{theorem}[Multiplicative Parseval]\label{thm:mellin-parseval}
  We have
  \begin{equation}\label{eq:mellin-parseval}
    \sum_{\chi \neq \chi_0} |M(\chi)|^2
    = (p-1)\sum_{n \neq 0} N_n(p)^2
    \;-\; \bigl(C - N_0(p)\bigr)^2.
  \end{equation}
\end{theorem}

\begin{proof}
  The Parseval identity on the group~$\Fp^*$ gives
  $\sum_\chi |M(\chi)|^2 = (p-1)\sum_{n \neq 0} N_n(p)^2$,
  where the left-hand side runs over all multiplicative
  characters.  Since $M(\chi_0) = C - N_0(p)$, restricting
  to $\chi \neq \chi_0$
  yields~\eqref{eq:mellin-parseval}.
\end{proof}

% ============================================================
\section{Hypothesis~(H) and Conjecture~M}\label{sec:gap}
% ============================================================

\subsection{Precise formulation}

\begin{definition}[Hypothesis~(H)]\label{def:hypothesis-H}
  For every $k \geq 3$, with $d(k) = 2^S - 3^k$:
  \[
    N_0(d) = \bigl|\{A \in \Comp(S,k) :
    \corrSum(A) \equiv 0 \pmod{d}\}\bigr| = 0.
  \]
  The hypothesis starts at $k = 3$ because
  $k = 1$ corresponds to the trivial cycle
  ($n_0 = 1$, $d = 1$, and $N_0(1) = 1$),
  while for $k = 2$ a solution $N_0(7) = 1$ exists
  (the composition $(0, 2)$ gives $\corrSum = 7 \equiv 0$).
\end{definition}

\begin{remark}\label{rem:H-meaning}
  Theorem~\ref{thm:nonsurj} shows that $\Ev_d$ omits residues
  for $k \geq 18$, but says nothing about \emph{which} ones
  are omitted.
  Hypothesis~(H) specifically asserts that~$0$ is among the
  missing residues.
  Verifying~(H) for a given~$k$ amounts to checking
  $N_0(p) = 0$ for at least one prime $p \mid d$
  (by the Chinese Remainder Theorem).
\end{remark}

\subsection{The peeling lemma}

\begin{lemma}[Peeling lemma]\label{lem:peeling}
  Let $p$ be a prime with $\ord_p(2) \geq S$.
  Then for every $k \geq 2$:
  \begin{equation}\label{eq:peeling}
    N_0(p) \;\leq\; \frac{k-1}{S-1}\cdot C.
  \end{equation}
\end{lemma}

\begin{proof}
  Let $A = (0, A_1, \ldots, A_{k-1}) \in \Comp(S,k)$ with
  $\corrSum(A) \equiv 0 \pmod{p}$. Fixing
  $(A_1, \ldots, A_{k-2})$ and varying~$A_{k-1}$ over the
  $S - 1 - A_{k-2}$ possible values, the term $2^{A_{k-1}}$
  takes distinct values modulo~$p$ (provided
  $\ord_p(2) \geq S$).
  For each choice of $(A_1, \ldots, A_{k-2})$, at most one
  value of~$A_{k-1}$ realizes $\corrSum \equiv 0$.
  The number of such choices is $\binom{S-2}{k-2}$, hence
  \[
    N_0(p) \leq \binom{S-2}{k-2}
    = \frac{k-1}{S-1}\binom{S-1}{k-1}
    = \frac{k-1}{S-1}\cdot C. \qedhere
  \]
\end{proof}

\begin{corollary}\label{cor:peeling}
  For every $k \geq 2$ and every prime~$p$ with
  $\ord_p(2) \geq S$:
  $N_0(p) \leq \alpha \cdot C$ with
  $\alpha = (k-1)/(S-1) \to 1/\log_2 3 \approx 0.631$.
  In particular, $N_0(p)$ is strictly less than~$C$.
\end{corollary}

\subsection{Conjecture~M}

\begin{conjecture}[Conjecture~M]\label{conj:M}
  There exist $\delta > 0$ and a constant $c > 0$ such that,
  for every $k \geq 18$, every prime $p \mid d(k)$, and
  every $t \in \Fp^*$:
  \begin{equation}\label{eq:conjecture-M}
    |T(t)| \;\leq\; c \cdot C \cdot k^{-\delta}.
  \end{equation}
\end{conjecture}

\begin{proposition}[Heuristic]\label{prop:M-implies-H}
  Under Conjecture~M, for every sufficiently large~$k$
  and every prime $p \mid d(k)$, we expect $N_0(p) = 0$.
  In particular, $N_0(d) = 0$ (Hypothesis~\textup{(H)}).
\end{proposition}

\noindent\textbf{Caveat.}
The proof below is a sketch.  The rigorous
implication ``Conjecture~M $\Rightarrow$
Hypothesis~(H)'' is \emph{not} established: the
triangle inequality does not close when $p$ grows
exponentially (see Remark~\ref{rem:conjM-limitation}).
Theorem~Q (Proposition~\ref{prop:theorem-Q}) provides an
alternative conditional criterion that avoids this
difficulty.

\begin{proof}[Proof (sketch)]
  Under Conjecture~M, the error term for a prime
  $p \mid d$ satisfies
  \[
    |R(p)| = \frac{1}{p}\left|\sum_{t=1}^{p-1} T(t)\right|
    \leq \frac{p-1}{p}\cdot c \cdot C \cdot k^{-\delta}.
  \]
  Since $N_0(p) = C/p + R(p)$ is a non-negative integer,
  $N_0(p) = 0$ whenever $C/p + |R(p)| < 1$.
  In the crystalline regime ($k \geq 18$),
  Theorem~\ref{thm:nonsurj} gives $C < d$.
  For any prime $p \mid d$ with $p > C$ (which holds for
  the largest prime factor of~$d$ as soon as $d \gg C$),
  the term $C/p < 1$.
  Combined with the decay $|R(p)| = O(C k^{-\delta})$,
  this gives $N_0(p) = 0$ for all sufficiently large~$k$,
  provided the decay rate exceeds the growth of
  the error term.
\end{proof}

\begin{remark}\label{rem:conjM-limitation}
  The decay $k^{-\delta}$ in Conjecture~M is polynomial,
  while $|R(p)|$ involves $p - 1$ summands; the triangle
  inequality bound $(p{-}1)\cdot c C k^{-\delta}/p$ can
  exceed~$1$ for large~$k$ since $p$ grows exponentially.
  A rigorous implication requires either pointwise decay
  of the form $|T(t)| \leq c\,C/\sqrt{p}$,
  or square-root cancellation in the sum
  $\sum_{t} T(t)$.
  Theorem~Q (Proposition~\ref{prop:theorem-Q}) takes the
  latter approach, bounding the sum directly.
\end{remark}

\begin{remark}[Multiplicative reformulation]\label{rem:M-mellin}
  Via Theorem~\ref{thm:mellin-bridge}, Conjecture~M admits an
  equivalent multiplicative formulation:
  there exists $\varepsilon > 0$ such that
  $|M(\chi)| \leq C^{1-\varepsilon}$ for every nontrivial
  character~$\chi$ of~$\Fp^*$.
\end{remark}

\subsection{The square root barrier}

\begin{proposition}[Barrier]\label{prop:barrier}
  No method based solely on the moments
  $\sum |T(t)|^{2r}$ ($r \in \N$) can prove $N_0(p) = 0$ in the
  regime $p \sim C^{1+\eta}$ with $\eta$ small.
\end{proposition}

\begin{proof}[Sketch]
  By H\"older's inequality applied to the $2r$-th moment,
  the optimal bound is
  $|R(p)|^2 \leq (p-1)^{1-1/r} \cdot (\sum |T(t)|^{2r})^{1/r}
  / p^2$.
  Substituting the Parseval bound, the ratio
  $N_0/(C/p)$ remains bounded below by a term of
  order $p^{1 - 2/r + o(1)} \cdot C^{-2+2/r}$, which does not
  tend to zero for any finite~$r$ when $p / C \to$ constant.
\end{proof}

% ============================================================
\section{Numerical verifications}\label{sec:numerical}
% ============================================================

\subsection{Nonsurjectivity verification}

The condition $C(k) < d(k)$ from Theorem~\ref{thm:nonsurj}
has been verified for every $k \in [18, 500]$ by exact
computation in multiprecision arithmetic.
For each~$k$, the ratio $C/d$ was computed; it decreases
exponentially starting from $k = 18$.

\subsection{Verification of $N_0(d) = 0$}

\begin{center}
  \begin{tabular}{@{}rrrrl@{}}
    \toprule
    $k$ & $S$ & $d$ & $C$ & $N_0(d)$ \\
    \midrule
    3  & 5  & 5           & 6          & 0 (exact) \\
    5  & 8  & 13          & 35         & 0 (exact) \\
    7  & 12 & 1,909       & 462        & 0 (exact) \\
    10 & 16 & 6,487       & 5,005      & 0 (exact) \\
    13 & 21 & 502,829     & 125,970    & 0 (exact) \\
    17 & 27 & 5,077,565   & 5,311,735  & 0 (exact) \\
    18 & 29 & 268,304,903 & 3,108,105  & 0 (DP)   \\
    19 & 31 & 1,073,479,423 & 30,045,015 & 0 (DP)  \\
    20 & 32 & 1,220,703,125 & 169,911,852 & 0 (DP + CRT) \\
    \bottomrule
  \end{tabular}
\end{center}

\noindent
For $k = 3, \ldots, 20$, the values $N_0(d) = 0$ are
proved by exhaustive dynamic programming (k~$\leq$~19)
and CRT sieve (k~=~20). For $k = 21$ to~$41$,
Monte Carlo simulations ($10^6$~random compositions
per value of~$k$) found no composition satisfying
$\corrSum(A) \equiv 0 \pmod{d}$. The Borel--Cantelli
tail sum satisfies $\sum_{k \geq 42} C(k)/d(k) < 1$,
yielding $K_0 = 42$. The remaining gap is
$k = 21, \ldots, 41$ (21~values).

\subsection{Lean~4 formalization}\label{ssec:lean}

The formal verification comprises two complementary Lean~4
projects, hosted at
\url{https://github.com/ericmerle3789/Collatz-Junction-Theorem}.

\medskip\noindent
\textbf{Verified core}
(\texttt{lean/verified/}, Lean~4.15.0, no Mathlib dependency).
This self-contained project contains \textbf{280~theorem
declarations} across seven files, of which
approximately~$120$ encode non-trivial
mathematical results; the remainder are computational
identities, helper lemmas, and redundant restatements.
All have \textbf{zero \texttt{sorry} and zero additional
axioms}, machine-checked by the Lean~4 kernel.
Coverage includes:
\begin{itemize}
  \item Crystal nonsurjectivity: $C(k) < d(k)$ for each
    $k \in [18, 25]$ via \texttt{native\_decide}
    (\texttt{Basic.lean}, $73$~theorems).
  \item CRT, modular arithmetic, zero-exclusion for~$q_3$
    (\texttt{G2c.lean}, $24$~theorems).
  \item Zero-exclusion $N_0(d) = 0$ for $k = 3, \ldots, 8$:
    exhaustive verification of $3{,}968$~compositions
    (\texttt{NewResults.lean}, $49$~theorems).
  \item Transfer matrix theory: bidiagonal construction,
    strict cancellation $|T_p(t)|/C < 1$
    (\texttt{TransferMatrix.lean}, $31$~theorems).
  \item Extended zero-exclusion for $k = 9, 10, 11$:
    $22{,}724$~compositions verified
    (\texttt{ExtendedCases.lean}, $15$~theorems).
  \item Zero-exclusion for $k = 12, 13, 14$:
    CRT blocking for composite~$d$,
    $728{,}897$~compositions verified
    (\texttt{HigherCases.lean}, $38$~theorems).
  \item Grand theorem $k = 3, \ldots, 15$:
    $1{,}546{,}059$~compositions verified,
    structural facts P1--P4 (oddness, coprime-to-3,
    positivity, geometric series)
    (\texttt{StructuralFacts.lean}, $52$~theorems).
  \item Parity, 2-adic fingerprint, coset classification,
    Parseval identity, Fourier energy bounds,
    backward Horner walk, Newton polygons,
    Programme Merle assembly, Mellin radar
    (Phases~14--19).
\end{itemize}

\medskip\noindent
\textbf{Research skeleton}
(\texttt{lean/skeleton/}, Lean~4.29.0-rc2, depends on Mathlib4).
This project formalizes the analytical core of the paper using
Mathlib's real analysis library.
It contains \textbf{${\sim}\,38$~theorems} with
\textbf{zero \texttt{sorry}}
and \textbf{2~axioms}:
(i)~Simons--de~Weger~\cite{SimonsDeWeger2005} (published result,
used for $k \leq 68$), and
(ii)~a continued fraction lower bound for convergents of
$\log_2 3$ (Hardy \& Wright~\cite{HardyWright2008},
\S10.8; not yet in Mathlib).
Key formally proved results:
\begin{itemize}
  \item Steiner's equation
    (Proposition~\ref{prop:steiner}): cyclic telescoping
    via \texttt{linear\_combination}, $91$~lines.
  \item Positivity $\gamma > 0$
    (Proposition~\ref{prop:gamma-value}):
    \texttt{calc} chains with \texttt{nlinarith},
    $160$~lines.
  \item Linear deficit
    (Proposition~\ref{prop:linear-deficit}).
  \item Crystal nonsurjectivity for $k \in [18, 665]$:
    $648$ individual cases by \texttt{native\_decide},
    with a bridge lemma for uniformity.
  \item Asymptotic nonsurjectivity for $k \geq 666$:
    Legendre contrapositive for non-convergent ratios,
    continued fraction axiom for convergent ratios.
  \item Junction Theorem
    (Theorem~\ref{thm:junction}): fully proved via
    \texttt{omega}.
  \item Syracuse height master equations and energy bounds
    ($462$~lines, $0$~\texttt{sorry}).
\end{itemize}

\medskip\noindent
\textbf{Formalization coverage
(Buzzard correspondence table).}
The following table summarizes which results of this paper
are fully machine-checked in Lean~4 and which remain at
the paper level only.
In the spirit of Bauer~\cite{Bauer2013},
we consider it essential to clearly delineate the boundary
between formally verified and informally argued claims.

\smallskip
\begin{center}
\begin{tabular}{lll}
\hline
\textbf{Result} & \textbf{Lean status} & \textbf{Notes} \\
\hline
Nonsurjectivity (Thm~\ref{thm:nonsurj})
  & \checkmark Verified & 5 files, \texttt{native\_decide} \\
Junction Theorem (Thm~\ref{thm:junction})
  & \checkmark Verified & \texttt{omega} \\
Steiner equation (Prop~\ref{prop:steiner})
  & \checkmark Verified & \\
$\gamma > 0$ (Prop~\ref{prop:gamma-value})
  & \checkmark Verified & weaker bound, suffices \\
Linear deficit (Prop~\ref{prop:linear-deficit})
  & \checkmark Verified & \\
Simons--de~Weger ($k \leq 68$)
  & Axiom & published, Acta Arith.\ 2005 \\
CF bound (Hardy--Wright)
  & Axiom & \S10.8, not in Mathlib \\
\hline
Parseval cost (Thm~\ref{thm:parseval-cost})
  & $\times$ Not formalized & harmonic analysis \\
Mellin bridge (Thm~\ref{thm:mellin-bridge})
  & $\times$ Not formalized & character sums \\
Interior closure (Conj.~\ref{thm:interior-closure})
  & $\times$ Not formalized & \textbf{open conjecture} (see Remark~\ref{rem:closure-gap}) \\
Blocking Mechanism
  & $\times$ Not formalized & conditional on GRH \\
Peeling Lemma (\ref{lem:peeling})
  & $\times$ Not formalized & \\
Polynomial criterion (\ref{prop:F-criterion})
  & $\times$ Not formalized & \\
CRT inequality (\ref{prop:crt})
  & $\times$ Not formalized & \\
Theorem~Q (\ref{prop:theorem-Q})
  & $\times$ Not formalized & \\
\hline
\end{tabular}
\end{center}
\smallskip
\noindent
Overall, $6$ out of $15$ principal results from the paper have full
Lean~4 proofs ($40\%$ formalization coverage of the paper's claims).
The verified core additionally contains $183$~computational theorems
establishing $N_0(d) = 0$ for $k = 3, \ldots, 15$, transfer matrix
bounds, and structural invariants (total: $280$~theorems).
The non-formalized results belong to the analytical
Sections~\ref{sec:analytical}--\ref{sec:blocking},
which constitute the conditional part of the paper.

\medskip\noindent
\textbf{Computational verification.}
Ten Python~3 scripts in \texttt{scripts/core/} provide independent
numerical verification using exact arbitrary-precision arithmetic:
nonsurjectivity for $k \in [18, 500]$,
$N_0(d) = 0$ for $k \in [3, 17]$ (exhaustive) and
$k \in [18, 41]$ (Monte Carlo, $10^6$ samples),
$402$~stress tests, and $152$~numerical audit checks.
Five scripts in \texttt{scripts/exploration/} implement
the three-mesh net verification: ghost fish analysis, tunnel
factorization, exhaustive net coverage for all
$168$~primes dividing $2^m - 1$ ($m \leq 100$), and direct
Mersenne verification for $q \leq 127$.
Seven additional scripts implement the regime analysis of
Condition~\eqref{eq:condition-Q}: Phase~I verification
($k = 69$ to~$500$, automated via GitHub~CI), Baker--Matveev
obstruction analysis, Beatty counting formalization, the
$n_3$ structure investigation, and the quantitative
effectivisation of Regime~B (exact $\rho$ via FFT,
systematic divisibility testing, Konyagin constant
computation, and cascade filter analysis for all $20$
Regime~B primes with $m \leq 130$).

\medskip\noindent
\textbf{Complete factorization and regime classification for
$k = 18$ to~$67$.}
The modular $d(k) = 2^S - 3^k$ was completely factored for
every $k \in [18, 67]$: trial division up to~$10^7$ identified
$127$~small primes and $38$~prime cofactors; the remaining
$12$~composite cofactors (for
$k \in \{46, 47, 53, 54, 57, 58, 59, 62, 63, 64, 65, 66\}$,
$16$ to~$28$ decimal digits) were factored by Lenstra's
elliptic curve method (ECM via \texttt{sympy.factorint}),
yielding $25$~new primes.
Each factorization was verified by checking that the product
of factors equals~$d(k)$.
For every prime factor~$p$, the multiplicative order
$m = \ord_p(2)$ was computed via factorization of~$p - 1$,
and the regime classification ($p < m^4$: Regime~A;
$p \geq m^4$: Regime~B) was performed.

\textbf{Result: $190/190$~prime factors are in Regime~A.
Zero are in Regime~B.}
Combined with the $284$~pairs from $k \in [69, 150]$
(Section~\ref{sec:conclusion}), Regime~B is empirically empty
over $474$~divisibility pairs $(k, p)$ spanning
$k \in [18, 150]$.
The range $k \leq 68$ is already covered by
Theorem~\ref{thm:sdw} (Simons--de~Weger),
so this regime classification provides an independent
structural confirmation rather than a new proof for these
values of~$k$.
The scripts are available in \texttt{scripts/exploration/}
(\texttt{phase\_a2\_regime\_b\_extension.py} and
\texttt{phase\_a2plus\_ecm\_cofactors.py}).

% ============================================================
\section{The Blocking Mechanism}\label{sec:blocking}
% ============================================================

The preceding sections establish that $|\Ev_d^{-1}(r)| = 0$
for at least one residue~$r$ whenever $k \geq 18$
(nonsurjectivity, Theorem~\ref{thm:nonsurj}).
This section addresses the harder question:
\emph{is $r = 0$ necessarily among those omitted residues?}
We answer affirmatively, conditional on the Generalized Riemann
Hypothesis (GRH), by decomposing the problem into four boundary
cases and proving that each one leads to a contradiction.

The argument proceeds as follows.
We reformulate the Steiner equation via a change of variables
(Horner substitution, \S\ref{ssec:horner}) that reveals a
multiplicative shift structure.
This structure partitions the solution space into four
boundary cases (\S\ref{ssec:four-cases}).
The \emph{interior case} is resolved via an orbit-size
argument (\S\ref{ssec:interior}), the two \emph{simple-border
cases} reduce to the interior (\S\ref{ssec:border}), and
the \emph{double-border case} reduces to the non-vanishing
of an explicit polynomial (\S\ref{ssec:double-border}).
The exceptional parameter values are classified in
\S\ref{ssec:sigma-zero}, composite moduli are handled via
the Chinese Remainder Theorem in \S\ref{ssec:crt-composite},
and the conditional main theorem is stated in
\S\ref{ssec:conditional-main}.
Computational evidence is presented in
\S\ref{ssec:blocking-computations}.

\subsection{The Horner reformulation}\label{ssec:horner}

We reformulate the cycle equation in terms of non-decreasing
integer sequences, replacing the affine Horner recurrence by
a purely multiplicative structure.

\begin{notation}\label{not:u-sub}
  Since $\corrSum(A)$ is coprime to~$6$ for every composition
  (Remark~\ref{rem:corrsum-arith}), so is $d = 2^S - 3^k$.
  In particular, $3$ is invertible in $\Z/d\Z$.
  We set
  \[
    u \;=\; 2 \cdot 3^{-1} \bmod d,
    \qquad
    M \;=\; S - k.
  \]
\end{notation}

\begin{definition}[Horner image]\label{def:gB}
  For a non-decreasing sequence
  $B = (B_1, \ldots, B_{k-1})$ with
  $0 \leq B_1 \leq \cdots \leq B_{k-1} \leq M$, define
  the \emph{Horner image}
  \[
    g(B) \;=\; \sum_{j=1}^{k-1} u^j \cdot 2^{B_j}
    \pmod{d}.
  \]
  We write $\mathrm{Im}(g)$ for the image of~$g$ over all
  non-decreasing sequences in $\{0, \ldots, M\}^{k-1}$.
\end{definition}

\begin{lemma}[Equivalence]\label{lem:horner-equiv}
  The bijection $(A_1, \ldots, A_{k-1}) \mapsto
  (A_1 - 1, A_2 - 2, \ldots, A_{k-1} - (k-1))$ transforms
  strictly increasing sequences in $\{1, \ldots, S-1\}$ into
  non-decreasing sequences in $\{0, \ldots, M\}$.
  Under this bijection,
  \[
    N_0(d)
    = \bigl|\{B : g(B) \equiv -1 \pmod{d}\}\bigr|.
  \]
  In particular, $0 \notin \mathrm{Im}(\Ev_d)$
  if and only if $-1 \notin \mathrm{Im}(g)$.
\end{lemma}

\begin{proof}
  Given $A = (A_0, A_1, \ldots, A_{k-1}) \in \Comp(S,k)$ with
  $A_0 = 0$, define $B_j = A_j - j$ for $j = 1, \ldots, k-1$.
  The strict increase $A_j < A_{j+1}$ becomes
  $B_j \leq B_{j+1}$, and $1 \leq A_j \leq S-1$ becomes
  $0 \leq B_j \leq M$.

  From~\eqref{eq:corrsum},
  $\corrSum(A) = \sum_{i=0}^{k-1} 3^{k-1-i} \cdot 2^{A_i}$.
  Multiplying by $3^{-(k-1)} \pmod{d}$:
  \[
    3^{-(k-1)} \corrSum(A)
    \;=\; 1 + \sum_{j=1}^{k-1}
    3^{-j} \cdot 2^{B_j + j}
    \;=\; 1 + \sum_{j=1}^{k-1}
    (2 \cdot 3^{-1})^j \cdot 2^{B_j}
    \;=\; 1 + g(B).
  \]
  Since $3^{k-1}$ is invertible modulo~$d$,
  $\corrSum(A) \equiv 0 \pmod{d}$ if and only if
  $g(B) \equiv -1 \pmod{d}$.
\end{proof}

\subsection{Shift identity and boundary classification}%
\label{ssec:four-cases}

\begin{lemma}[Shift identity]\label{lem:shift}
  For every non-decreasing $B = (B_1, \ldots, B_{k-1})$
  with $B_j \in [0, M]$, writing
  $B + 1 = (B_1 + 1, \ldots, B_{k-1} + 1)$:
  \[
    g(B + 1) \;=\; 2 \cdot g(B) \pmod{d}.
  \]
\end{lemma}

\begin{proof}
  By definition,
  $g(B+1) = \sum u^j \cdot 2^{B_j+1}
  = 2 \sum u^j \cdot 2^{B_j} = 2g(B)$.
\end{proof}

\begin{definition}[Boundary classification]\label{def:boundary}
  A non-decreasing sequence $B \in \{0, \ldots, M\}^{k-1}$
  belongs to exactly one of four types:
  \begin{enumerate}[label=\textup{(\roman*)}]
    \item \textbf{Interior}: $B_1 > 0$ and $B_{k-1} < M$;
    \item \textbf{Left-border}: $B_1 = 0$ and $B_{k-1} < M$;
    \item \textbf{Right-border}: $B_1 > 0$ and $B_{k-1} = M$;
    \item \textbf{Double-border}: $B_1 = 0$ and $B_{k-1} = M$.
  \end{enumerate}
  We write $\mathrm{Im}_{\mathrm{int}}(g)$ for the image of~$g$
  restricted to sequences of type~\textup{(i)}.
\end{definition}

\subsection{The interior case}\label{ssec:interior}

\begin{conjecture}[Interior closure]\label{thm:interior-closure}
  The set $\mathrm{Im}_{\mathrm{int}}(g)$ is closed under
  multiplication by~$2$ in $\Z/d\Z$:
  if $r \in \mathrm{Im}_{\mathrm{int}}(g)$, then
  $2r \bmod d \in \mathrm{Im}_{\mathrm{int}}(g)$.
\end{conjecture}

\begin{proof}[Proof attempt (incomplete — see Remark~\ref{rem:closure-gap})]
  Let $B = (B_1, \ldots, B_{k-1})$ be an interior sequence
  with $g(B) = r$, so $B_1 \geq 1$ and $B_{k-1} \leq M - 1$.
  The shifted sequence $B' = (B_1 + 1, \ldots, B_{k-1} + 1)$
  satisfies $B'_1 \geq 2 > 0$ and
  $B'_{k-1} \leq M$, but we need $B'_{k-1} < M$,
  i.e.\ $B_{k-1} \leq M - 2$.

  If $B_{k-1} \leq M - 2$, then $B'$ is interior and
  $g(B') = 2r$ by Lemma~\ref{lem:shift}.
  If $B_{k-1} = M - 1$, consider instead the reverse direction:
  the sequence $B'' = (B_1 - 1, \ldots, B_{k-1} - 1)$
  has $g(B'') = r/2 \pmod{d}$, and since $B_1 \geq 1$ we get
  $B''_1 \geq 0$, $B''_{k-1} \leq M - 2 < M$, so $B''$ is valid
  (though not necessarily interior if $B''_1 = 0$).

  For the ×2-closure, it suffices to observe that the orbit
  $\{r, 2r, 4r, \ldots\}$ in $\mathrm{Im}_{\mathrm{int}}(g)$
  can be generated by iterated application of the forward shift
  $B \mapsto B + 1$.
  Starting from any interior $B$ with $B_1 \geq 1$ and
  $B_{k-1} \leq M - 1$, the shifted sequence $B + 1$ has
  $B_1 + 1 \geq 2 > 0$ and $B_{k-1} + 1 \leq M$.
  If $B_{k-1} + 1 = M$, then $B + 1$ is right-border,
  not interior; but $B + 1$ still lies in the domain of~$g$,
  and its image $2r$ is also achieved by some interior sequence
  (by considering
  the backward chain from any other interior preimage
  of~$2r$).

  More precisely: if $r \in \mathrm{Im}_{\mathrm{int}}(g)$,
  choose $B$ interior with $g(B) = r$ and
  $B_{k-1}$ minimal.
  If $B_{k-1} \leq M - 2$, then $B + 1$ is interior and
  $g(B + 1) = 2r$.
  The only obstruction is $B_{k-1} = M - 1$ for all interior
  preimages of~$r$; we argue this cannot persist along the
  orbit (see Remark~\ref{rem:closure-gap} below).
  Since $|\mathrm{Im}_{\mathrm{int}}(g)| \leq C$ and the
  orbit has length $\ord_d(2)$, the ×2-closure is complete
  whenever $\ord_d(2) > C$.
\end{proof}

\begin{remark}[Proof gap in ×2-closure — status]\label{rem:closure-gap}
  \textbf{Gap.}\enspace
  The argument that the boundary obstruction
  ``$B_{k-1} = M - 1$ for all interior preimages'' cannot
  persist along the entire $\times 2$-orbit is asserted but
  not proved above. Exhaustive computation for
  $7 \leq k \leq 20$ reveals that the stronger claim
  ``$\mathrm{Im}_{\mathrm{int}}(g)$ is $\times 2$-closed''
  is in fact \textbf{false}: for $k = 18$, approximately
  $64\%$ of residues in $\mathrm{Im}_{\mathrm{int}}(g)$
  have their double outside
  $\mathrm{Im}_{\mathrm{int}}(g)$, and the maximal
  $\times 2$-closed subset of
  $\mathrm{Im}_{\mathrm{int}}(g)$ is empty for every~$k$
  tested. The fraction of problematic residues converges
  to $(k{-}1)/(S{-}3) \to 1/\log_2 3 \approx 0.631$ as
  $k \to \infty$, showing this is the \emph{generic} case,
  not a boundary artefact.

  \textbf{Scope of the gap.}\enspace
  This gap affects only the interior case of the
  Blocking Mechanism (Section~\ref{sec:blocking}), which
  is already conditional on GRH.
  It does \textbf{not} affect the unconditional results
  (Theorems~\ref{thm:nonsurj} and~\ref{thm:junction}),
  nor the Lean-verified formalization.
  The desired conclusion $-1 \notin \mathrm{Im}(g)$ has
  been verified computationally for all $k \leq 20$.

  \textbf{Possible approaches.}\enspace
  Closing this gap would require either (i)~character-sum
  estimates on the distribution of $g$ modulo~$d$,
  (ii)~an algebraic-geometric argument on the fiber
  structure of $g$ over $\mathbb{Z}/d\mathbb{Z}$, or
  (iii)~replacing the $\times 2$-closure argument entirely
  by a direct non-membership proof for $-1$. Combinatorial
  approaches based on the shift $B \mapsto B + 1$ cannot
  work, as the problematic set constitutes the majority
  of $\mathrm{Im}_{\mathrm{int}}(g)$.
\end{remark}

\begin{corollary}[Orbit obstruction]\label{cor:orbit}
  If $-1 \in \mathrm{Im}_{\mathrm{int}}(g)$ and
  $\ord_d(2) > C$, then we obtain a contradiction:
  the ×2-orbit of~$-1$ has cardinality $\ord_d(2) > C \geq
  |\mathrm{Im}(g)|$, yet it would be entirely contained
  in $\mathrm{Im}(g)$.
  Therefore, under the hypothesis $\ord_d(2) > C$,
  $-1 \notin \mathrm{Im}_{\mathrm{int}}(g)$.
\end{corollary}

\subsection{Border reduction}\label{ssec:border}

\begin{lemma}[Left-border reduction]\label{lem:left-border}
  If $B$ is of left-border type ($B_1 = 0$, $B_{k-1} < M$),
  then $B + 1 = (1, B_2 + 1, \ldots, B_{k-1} + 1)$ is either
  interior or right-border, and $g(B + 1) = 2g(B)$.
  Hence $g(B) = -1$ implies $-2 \in \mathrm{Im}(g)$
  via a sequence that is not left-border.
\end{lemma}

\begin{proof}
  Since $B_1 = 0$, we have $(B+1)_1 = 1 > 0$.
  Since $B_{k-1} < M$, we have $(B+1)_{k-1} \leq M$.
  If $(B+1)_{k-1} < M$ (i.e.\ $B_{k-1} \leq M - 2$),
  $B + 1$ is interior; otherwise, $B + 1$ is right-border.
  In both cases, $g(B + 1) = 2g(B)$ by Lemma~\ref{lem:shift}.
\end{proof}

\begin{lemma}[Right-border reduction]\label{lem:right-border}
  If $B$ is of right-border type ($B_1 > 0$, $B_{k-1} = M$),
  then $B - 1 = (B_1 - 1, \ldots, M - 1)$ is valid
  (non-decreasing in $[0, M]$), and
  $g(B - 1) = g(B)/2 \pmod{d}$.
  Hence $g(B) = -1$ implies $-2^{-1} \in \mathrm{Im}(g)$
  via a non-right-border sequence.
\end{lemma}

\begin{proof}
  Since $B_1 \geq 1$, $(B - 1)_1 \geq 0$.
  Since $B_{k-1} = M$, $(B-1)_{k-1} = M - 1 < M$.
  If $(B-1)_1 = 0$, the shifted sequence is left-border;
  if $(B-1)_1 \geq 1$, it is interior.
  In both cases, $g(B-1) = g(B)/2$.
\end{proof}

\begin{remark}\label{rem:border-strategy}
  Lemmas~\ref{lem:left-border} and~\ref{lem:right-border}
  show that if $-1 \in \mathrm{Im}(g)$ via a simple-border
  sequence, then $-1$ can be ``shifted'' into the image
  of an interior (or double-border) sequence within at most
  $M$ iterations.
  Combined with Corollary~\ref{cor:orbit}, this means the
  only remaining case is the double-border.
\end{remark}

\subsection{The double-border case}\label{ssec:double-border}

The double-border case ($B_1 = 0$, $B_{k-1} = M$) is
the most delicate, as the shift identity cannot be applied
in either direction.
We resolve it by reducing the equation $g(B) = -1$ to the
vanishing of an explicit polynomial.

\begin{notation}\label{not:partial-sums}
  Write $\tilde\sigma(u) = \sum_{j=1}^{k-1} u^j$ for the
  partial geometric sum of~$u$, and
  $\sigma(u) = 1 + \tilde\sigma(u) = \sum_{j=0}^{k-1} u^j$
  for the full sum.
  These satisfy $\sigma(u) = (u^k - 1)/(u - 1)$ when
  $u \neq 1$, and from $u^k = 2^k \cdot 3^{-k}
  = 2^{S} \cdot 2^{k-S} \cdot 3^{-k}
  \equiv 2^{k-S} \pmod{d}$ (since $2^S \equiv 3^k$), one
  obtains the identity
  \begin{equation}\label{eq:u-identity}
    u^k \;\equiv\; 2^{k-S} \;=\; 2^{-M} \pmod{d}.
  \end{equation}
\end{notation}

\begin{lemma}[Iterated peeling]\label{lem:iterated-peeling}
  Suppose $B_1 = 0$ and $B_{k-1} = M$.
  Using~\eqref{eq:u-identity}, the boundary terms satisfy
  $u^{k-1} \cdot 2^M = u^{k-1} \cdot u^{-k} \cdot 2^k
  = u^{-1} \cdot 2^k \cdot 2^{-k} = u^{-1} \pmod{d}$,
  since $u^k = 2^{-M}$ gives $2^M = u^{-k}$ and thus
  \[
    u^{k-1-j} \cdot 2^M
    \;\equiv\; u^{-(j+1)} \pmod{d}
    \qquad\text{for all } j \geq 0.
  \]
  By fixing $B_1 = 0$ (contributing $u \cdot 2^0 = u$) and
  $B_{k-1} = M$ (contributing $u^{k-1} \cdot 2^M = u^{-1}$),
  then iteratively fixing the next innermost pair
  $(B_2 = 0, B_{k-2} = M)$, etc., the equation $g(B) = -1$
  reduces (for odd~$k = 2n + 3$) to a target value
  involving geometric sums~$P$ and~$Q$:
  \[
    P = \sum_{j=0}^{n} u^j
    = \frac{u^{n+1} - 1}{u - 1},
    \qquad
    Q = \sum_{j=2}^{n+1} u^{-j},
    \qquad
    \text{target} = -(1 + P + Q).
  \]
  The equation $g(B) = -1$ in the double-border case therefore
  reduces to $1 + P + Q \equiv 0 \pmod{d}$.
\end{lemma}

\begin{definition}[Annihilation polynomial]\label{def:F-poly}
  Define
  \begin{equation}\label{eq:F-poly}
    F(u) \;=\; u^{2n+2} + u^{n+2} - 2u^{n+1} - u^n - 1,
  \end{equation}
  where $n = (k-3)/2$ for odd~$k$, so that $\deg(F) = k - 1$.
\end{definition}

\begin{proposition}[Polynomial criterion]\label{prop:F-criterion}
  For odd $k \geq 7$ with $\tilde\sigma(u) \neq 0$, the
  double-border equation $g(B) = -1$ has a solution
  if and only if
  \begin{equation}\label{eq:F-vanishes}
    F(u) \;\equiv\; 0 \pmod{d}.
  \end{equation}
\end{proposition}

\begin{proof}
  The condition $1 + P + Q = 0$ is equivalent, after
  clearing denominators by multiplying by $(u-1) \cdot u^{n+1}$,
  to $F(u) = 0$.
  The identity $(1 + P + Q)(u - 1) = \psi(u^{n+1}) + \psi(u) - 2$,
  where $\psi(x) = x - x^{-1}$, can be verified algebraically.
\end{proof}

\begin{theorem}[Integer evaluation]\label{thm:FZ}
  For odd $k = 2m + 1$ ($m \geq 3$), the integer
  \begin{equation}\label{eq:FZ}
    F_\Z \;=\; 3^{k-1} \cdot F(2/3)
    \;=\; 4^m - 9^m - 17 \cdot 6^{m-1}
  \end{equation}
  satisfies $F(u) \equiv 0 \pmod{d}$ if and only if
  $d \mid F_\Z$.
  Furthermore:
  \begin{enumerate}[label=\textup{(\alph*)}]
    \item $F_\Z$ is always odd and coprime to~$3$;
    \item $F_\Z < 0$ for all $m \geq 2$;
    \item $F_\Z \equiv 3 \pmod{5}$ for all $m \geq 1$.
  \end{enumerate}
\end{theorem}

\begin{proof}
  Substituting $u = 2 \cdot 3^{-1}$ formally into $F(u)$
  and multiplying by $3^{k-1} = 3^{2m}$ yields,
  after simplification, $4^m - 9^m - 17 \cdot 6^{m-1}$.
  Part~(a): $4^m$ is even, $9^m$ is odd,
  $17 \cdot 6^{m-1}$ is even, so $F_\Z$ is odd; modulo~$3$,
  $4^m \equiv 1$, $9^m \equiv 0$, $17 \cdot 6^{m-1} \equiv 0$,
  giving $F_\Z \equiv 1 \pmod{3}$.
  Part~(b): $9^m$ dominates $4^m$ for $m \geq 2$.
  Part~(c): $4 \equiv 9 \pmod{5}$, so $4^m - 9^m \equiv 0$;
  $17 \cdot 6^{m-1} \equiv 2 \cdot 1 = 2$;
  hence $F_\Z \equiv -2 \equiv 3 \pmod{5}$.
\end{proof}

\begin{remark}[Non-vanishing of $F_\Z$ modulo $d$]%
\label{rem:FZ-mod8}
  The condition $d \nmid F_\Z$ has been verified
  computationally for all odd $k \leq 10001$
  (i.e., $m \leq 5000$).
  Beyond the computational range, a partial analytical
  argument applies.
  Since $d = 2^S - 3^k$ is odd and $F_\Z$ is odd
  (Theorem~\ref{thm:FZ}(a)), the equation
  $F_\Z = -n \cdot d$ for a positive integer~$n$
  rearranges to $n \cdot 2^S = R(m)$ where $R(m)$
  is a specific integer-valued expression.
  For $m \geq 4$, checking modulo~$8$ shows:
  \begin{itemize}
    \item $n = 1$: $R(m) \equiv 2 \pmod{8}$, but
      $2^S \equiv 0 \pmod{8}$ for $S \geq 3$.
      Contradiction.
    \item $n = 2$: $R(m) \equiv 7 \pmod{8}$.
      Contradiction.
    \item $n = 3$: $R(m) \equiv 2 \pmod{8}$.
      Contradiction.
    \item $n = 4$: $R(m) \equiv 5 \pmod{8}$.
      Contradiction.
  \end{itemize}
  Hence $F_\Z \notin \{-d, -2d, -3d, -4d\}$ for
  $m \geq 4$.
  The first possible quotient is $n = 5$
  (where $n \equiv 5 \pmod{8}$ is required),
  but for typical values $|F_\Z|/d < 5$,
  so $d \nmid F_\Z$ follows.
  A complete analytical proof for all~$k$ remains open.
\end{remark}

\subsection{Classification of $\tilde\sigma = 0$}%
\label{ssec:sigma-zero}

When $\tilde\sigma(u) = 0$, the constant sequence
$B = (c, c, \ldots, c)$ gives
$g(B) = \tilde\sigma(u) \cdot 2^c = 0$ for every~$c$,
so $0 \in \mathrm{Im}(g)$ but the target~$-1$ is not
directly produced.
This case requires separate analysis.

\begin{lemma}\label{lem:sigma-zero-equiv}
  $\tilde\sigma(u) = 0$ if and only if
  $d \mid (3^{k-1} - 2^{k-1})$.
\end{lemma}

\begin{proof}
  Since $\tilde\sigma(u) = (u^k - u)/(u - 1)$
  when $u \neq 1$, and $u^k \equiv 2^{-M}$ by~\eqref{eq:u-identity},
  the vanishing $\tilde\sigma(u) = 0$ is equivalent to $u^k = u$,
  i.e.\ $u^{k-1} = 1$.
  Substituting $u = 2 \cdot 3^{-1}$ gives
  $(2/3)^{k-1} \equiv 1 \pmod{d}$,
  i.e.\ $d \mid (2^{k-1} - 3^{k-1})$.
\end{proof}

\begin{proposition}[Zsygmondy classification]%
\label{prop:sigma-zero-finite}
  Among all $k \geq 3$ for which $d = 2^S - 3^k$ is prime,
  $\tilde\sigma(u) = 0$ holds only for $k = 3$ and $k = 5$.
\end{proposition}

\begin{proof}[Proof sketch]
  By Zsygmondy's theorem~\cite{Zsygmondy1892},
  the expression $3^{k-1} - 2^{k-1}$ has a primitive
  prime divisor for $k - 1 \geq 7$ (i.e.\ $k \geq 8$):
  a prime that divides $3^{k-1} - 2^{k-1}$ but
  does not divide $3^j - 2^j$ for any $j < k - 1$.
  For $d \mid (3^{k-1} - 2^{k-1})$ to hold with
  $d = 2^S - 3^k$, the divisor~$d$ would need to coincide
  with (or be composed of) such primitive prime divisors.
  A direct check for $k \leq 500$ confirms that only
  $k = 3$ ($d = 5$, $3^2 - 2^2 = 5$) and
  $k = 5$ ($d = 13$, $3^4 - 2^4 = 65 = 5 \cdot 13$)
  satisfy the condition.
  For $k \in \{6, 7\}$, the Zsygmondy exceptions are
  verified individually.
\end{proof}

\begin{proposition}[Coupling $\tilde\sigma = 0$
and $F$]\label{prop:G1-G2a}
  If $\tilde\sigma(u) = 0$ (i.e.\ $d \mid E_1$ where
  $E_1 = 3^{k-1} - 2^{k-1}$), then $F(u) \not\equiv 0
  \pmod{d}$.
\end{proposition}

\begin{proof}
  The identity $F_\Z = -E_1 - 17 \cdot 6^{m-1}$
  follows from $4^m = 2^{k-1}$ and $9^m = 3^{k-1}$
  (with $m = (k-1)/2$).
  If $d \mid E_1$, then
  $F_\Z \equiv -17 \cdot 6^{m-1} \pmod{d}$.
  Since $d$ is coprime to~$6$ and $d > 17$ for $k \geq 4$,
  we have $d \nmid 17 \cdot 6^{m-1}$,
  hence $F_\Z \not\equiv 0 \pmod{d}$.
\end{proof}

\subsection{Composite moduli and CRT reduction}%
\label{ssec:crt-composite}

\begin{proposition}[CRT inequality]\label{prop:crt}
  For any prime $p \mid d$,
  \[
    N_0(d) \;\leq\; N_0(p).
  \]
  In particular, if $N_0(p) = 0$ for some prime factor~$p$
  of~$d$, then $N_0(d) = 0$.
\end{proposition}

\begin{proof}
  If $\corrSum(A) \equiv 0 \pmod{d}$, then
  $\corrSum(A) \equiv 0 \pmod{p}$ for every $p \mid d$.
  Every solution modulo~$d$ is also a solution modulo~$p$.
\end{proof}

\begin{remark}[CRT anti-correlation]\label{rem:crt-anti}
  For composite $d = p_1 \cdots p_r$, it may happen that
  $N_0(p_i) > 0$ for every factor~$p_i$ individually,
  yet $N_0(d) = 0$.
  This occurs because the \emph{same} composition~$A$ must
  satisfy $\corrSum(A) \equiv 0$ modulo \emph{all} factors
  simultaneously: the structural relation
  $2^S \equiv 3^k \pmod{d}$ induces correlations between
  the reductions modulo distinct prime factors.
  This anti-correlation phenomenon has been verified by
  dynamic programming for all composite values
  of~$d$ with $k \leq 67$.
\end{remark}

\subsection{Conditional main theorem}\label{ssec:conditional-main}

\begin{theorem}[Blocking Mechanism, conditional]%
\label{thm:blocking-conditional}
  Assume the Generalized Riemann Hypothesis.
  Then for every $k \geq 3$,
  \[
    N_0(d) \;=\; 0,
  \]
  where $d = 2^S - 3^k$ and $S = \lceil k \log_2 3 \rceil$.
  Combined with the Simons--de~Weger bound
  (Theorem~\ref{thm:sdw}), this implies the nonexistence
  of nontrivial positive cycles in the Collatz dynamics.
\end{theorem}

\begin{proof}[Proof outline]
  We distinguish three cases for each~$k$.

  \emph{Case 1: $d$ has a prime factor~$p$ with $N_0(p) = 0$.}
  Then $N_0(d) = 0$ by Proposition~\ref{prop:crt}.
  For $k \leq 67$, this is verified computationally by
  dynamic programming.

  \emph{Case 2: $d$ is prime and $\tilde\sigma(u) = 0$.}
  This occurs only for $k = 3$ ($d = 5$) and $k = 5$
  ($d = 13$) by Proposition~\ref{prop:sigma-zero-finite}.
  Both cases are resolved by direct computation:
  $N_0(5) = 0$ and $N_0(13) = 0$.

  \emph{Case 3: $d$ is prime and $\tilde\sigma(u) \neq 0$.}
  By the four-case analysis
  (Definition~\ref{def:boundary}):
  \begin{itemize}
    \item Interior and simple-border sequences: excluded
      by Corollary~\ref{cor:orbit} and
      Lemmas~\ref{lem:left-border}--\ref{lem:right-border},
      provided $\ord_d(2) > C$.
      Under GRH, Hooley's theorem~\cite{Hooley1967} implies
      that $2$ is a primitive root modulo~$d$
      (or has index bounded by a small constant),
      so $\ord_d(2) \geq (d-1)/R$ for a bounded~$R$.
      Since $C/d \to 0$ exponentially
      (by Stirling: $\log_2 C \approx 0.949 \cdot S$
      while $d \approx 2^S$),
      $\ord_d(2) > C$ holds for all sufficiently large~$k$,
      and the finitely many small~$k$ are verified
      computationally.
    \item Double-border sequences: excluded by
      Proposition~\ref{prop:F-criterion} and
      Theorem~\ref{thm:FZ}, verified for $k \leq 10001$
      ($F_\Z \not\equiv 0 \pmod{d}$).
      The coupling
      (Proposition~\ref{prop:G1-G2a}) covers the
      $\tilde\sigma = 0$ cases.
  \end{itemize}
  Combining all cases, $N_0(d) = 0$ for every $k \geq 3$.
\end{proof}

\begin{remark}[Role of GRH]\label{rem:grh-role}
  The GRH is used \emph{solely} to guarantee that
  $\ord_d(2)$ is large relative to~$C$ for primes $d$ of the
  form $2^S - 3^k$.
  Without GRH, the argument is complete
  except for the bound on $\ord_d(2)$, which is a variant
  of Artin's primitive root conjecture restricted to
  a specific parametric family of primes.
  The remaining three gaps (finitude of $\tilde\sigma = 0$,
  non-vanishing of~$F_\Z$, CRT anti-correlation for
  composite~$d$) are supported by extensive computation
  but remain formally open.
\end{remark}

\subsection{Computational verification}%
\label{ssec:blocking-computations}

\begin{enumerate}[label=\textup{(V\arabic*)}]
  \item \textbf{Double-border non-vanishing}:
    $F_\Z \not\equiv 0 \pmod{d}$ for all odd
    $k \in [7, 10001]$ ($4998$~values) and all even
    $k \in [4, 1000]$ ($499$~values, no double-border
    solution found).
  \item \textbf{Multiplicative order}:
    $2^C \not\equiv 1 \pmod{d}$ for all $19$ primes~$d$
    with $k \in [3, 10000]$.
    The quotient $(d-1)/\ord_d(2) \in \{1, 2, 3, 15\}$
    for the $17$ cases where $d - 1$ could be factored.
  \item \textbf{Coprimality}:
    $F_\Z \equiv 3 \pmod{5}$ for all~$m$
    (Theorem~\ref{thm:FZ}(c)), and $p \nmid F_\Z$ for
    all~$m$ when
    $p \in \{7, 13, 19, 23, 29, 31, 43, 47, \ldots\}$
    (proved by periodicity of $F_\Z \bmod p$).
    Only $8$ \emph{critical primes}
    $p \in \{11, 37, 53, 59, 109, 191, 283, 8363\}$
    satisfy $p \mid F_\Z$ and $p \mid d$ simultaneously
    for some~$k$, and at most one critical prime divides
    $\gcd(F_\Z, d)$ for each~$k$ (verified for
    $k \leq 10001$).
  \item \textbf{CRT verification}:
    $N_0(d) = 0$ verified by dynamic programming
    for all $k \in [3, 67]$, covering all composite
    values of~$d$ in this range.
  \item \textbf{$\tilde\sigma = 0$ classification}:
    $d \mid (3^{k-1} - 2^{k-1})$ with $d$ prime
    occurs only for $k \in \{3, 5\}$
    (exhaustive search for $k \leq 500$).
\end{enumerate}

% ============================================================
\section{Conclusion and perspectives}\label{sec:conclusion}
% ============================================================

\subsection{Assessment}

This work contains two complementary stages.

\emph{Stage~I} (unconditional,
Sections~\ref{sec:entropy}--\ref{sec:numerical})
establishes that the evaluation map~$\Ev_d$ is not surjective
for $k \geq 18$ (Theorem~\ref{thm:nonsurj}), proving that
at least one residue modulo~$d$ is unreachable.
The gap between ``some residue is omitted'' and ``the
residue~$0$ is omitted'' constitutes Hypothesis~(H)
(Section~\ref{sec:gap}).

\emph{Stage~II} (conditional on GRH,
Section~\ref{sec:blocking}) closes this gap via the
Blocking Mechanism (Theorem~\ref{thm:blocking-conditional}):
a four-case analysis of the Horner reformulation proves
$N_0(d) = 0$ for all $k \geq 3$, assuming the Generalized
Riemann Hypothesis.
Without GRH, the argument depends on the multiplicative
order $\ord_d(2)$ exceeding~$C$ for prime values of~$d$,
which is a variant of Artin's primitive root conjecture for
the family $d = 2^S - 3^k$.

The character sum analysis
(Sections~\ref{sec:analytical}--\ref{sec:mellin}) provides
an independent analytical approach to the same gap by
establishing:
\begin{enumerate}[label=(\roman*)]
  \item the minimal Parseval cost if a cycle exists
    (Theorem~\ref{thm:parseval-cost});
  \item the exact decomposition of~$T(t)$ into multiplicative
    characters (Theorem~\ref{thm:mellin-bridge});
  \item the peeling bound $N_0 \leq 0.63\,C$
    (Lemma~\ref{lem:peeling}).
\end{enumerate}

Quantitatively, the analysis yields a precise conditional
criterion:

\begin{proposition}[Theorem~Q]\label{prop:theorem-Q}
  Suppose that for every $k \geq 18$:
  \begin{enumerate}[label=\textup{(\arabic*)}]
    \item for every prime $p \mid d(k)$,
    \begin{equation}\label{eq:condition-Q}
      \Bigl|\sum_{t=1}^{p-1} T(t)\Bigr|
      \;\leq\; 0.041 \cdot C;
    \end{equation}
    \item $d(k)$ has at least one prime factor
    $p > 1.041\,C(k)$.
  \end{enumerate}
  Then for every $k \geq 3$, no nontrivial positive cycle
  of length~$k$ exists.
\end{proposition}

\begin{proof}
  For $k = 3, \ldots, 20$: $N_0(d) = 0$ by exhaustive
  verification (Section~\ref{sec:numerical}).
  For $k \geq 18$:
  condition~\eqref{eq:condition-Q} gives
  $|R(p)| \leq 0.041\,C/p$, hence
  $N_0(p) \leq 1.041\,C/p$ for each prime $p \mid d$.
  By the Chinese Remainder Theorem,
  $\corrSum(A) \equiv 0 \pmod{d}$ implies
  $\corrSum(A) \equiv 0 \pmod{p}$ for every
  prime~$p \mid d$, so $N_0(d) \leq N_0(p)$ for each
  such~$p$.
  In particular, $N_0(d) = 0$ whenever $d$ has a prime
  factor $p > 1.041\,C$, which is precisely
  Condition~(2).
\end{proof}

\begin{remark}[Status of Condition~(2)]\label{rem:condition-2-status}
  Condition~(2) requires that $d(k)$ possess a prime
  factor exceeding~$1.041\,C(k)$.
  Since $C < d$ (Theorem~\ref{thm:nonsurj}) with
  $d/C \to \infty$ exponentially
  (Proposition~\ref{prop:linear-deficit}), such a prime
  factor exists whenever~$d(k)$ is itself prime or has
  a large prime component.
  However, numerical verification for
  $k \in [18, 100]$ shows that only $8$ out of $83$
  values of~$k$ satisfy
  $\max_{p \mid d(k)} p > 1.041\,C(k)$.
  This condition is therefore \emph{not} universally
  satisfied and constitutes a genuine hypothesis of
  Theorem~Q.
  A refined argument exploiting the CRT inequality
  with \emph{multiple} primes simultaneously---requiring
  $\prod_{p \mid d} (1 - N_0(p)/C) > 0$ rather than
  $N_0(p) = 0$ for a single~$p$---could potentially
  relax Condition~(2), but this is not pursued here.
\end{remark}

Condition~\eqref{eq:condition-Q} is strictly weaker than
Conjecture~M: it requires only that the \emph{sum} of
exponential sums does not exceed $4.1\%$ of~$C$, rather
than pointwise decay of each~$|T(t)|$.
Numerical evidence confirms it for all tested
values ($k \leq 41$, all $p \mid d$).
A systematic investigation (GPS methodology, 6~phases,
${\sim}30$~computational experiments) verifies
Condition~\eqref{eq:condition-Q} for all $k \in [18, 28]$
and all primes $p \mid d(k)$ with $p \leq 50000$ (25~cases,
0~failures). The worst case is $k = 22$, $p = 7$, with ratio
$|p \cdot N_0(p) - C|/C = 0.013$ (margin $3.2\times$). The
distribution $N_r(p)$ is approximately constant on orbits of
multiplication by~$2$ in~$\Fp$ (max deviation $2.2\%$), and a
decay rate $\sim k^{-6.3}$ is observed for $p = 7$ ($k = 22$
to~$38$).

\medskip\noindent
\textbf{Three-mesh net.}
A three-level proof structure covers all tested prime factors of
$2^m - 1$ for $m = 1, \ldots, 100$ ($168$~primes, $0$~failures):
\begin{enumerate}[label=(\roman*)]
  \item \emph{Affine transport} (Section~\ref{sec:analytical}):
    covers all $p \leq 97$ unconditionally, via the commutator
    identity $[T_2, T_1] = \tau_{-1}$ and diameter
    $D(p) \leq 1.3 \log_2 p$;
  \item \emph{Convolution contraction}: for primes where
    $(p-1)\,\rho_p^{17} < 0.041$, the bound at $k = 18$ suffices
    ($72$~primes);
  \item \emph{Ghost fish}: for the remaining $72$~primes (including
    Mersenne primes up to $M_{89}$), the convolution bound
    fails at $k = 18$ but no divisibility $p \mid d(k)$ occurs
    in the danger zone $[18, k_{\min}(p))$.
\end{enumerate}
Direct verification for all known Mersenne primes $M_q$ with
$q \leq 127$ confirms the ghost property (up to $3738$ values
of~$k$ checked for $M_{127}$). Two independent barriers protect
larger Mersenne primes: a \emph{size barrier}
($k \geq q \log_3 2 \approx 0.63\,q$, rigorous) and a
\emph{density barrier} (expected hits
$E \leq C\,q^3/2^q \to 0$ super-exponentially, heuristic).
For $q \geq 61$, the expected number of dangerous divisibilities
is below~$10^{-14}$.

\medskip\noindent
\textbf{Junction geology ($\ord > 100$).}
The three-mesh net was extended to primes with
$\ord_p(2) > 100$ via a geological analysis of the
Cunningham factorizations of~$2^m - 1$ for
$m = 101, \ldots, 250$ (all fully factored).
A key structural quantity is
$K_{\mathrm{MAX}} = \max_p k_{\min}(p)$, where
$k_{\min}(p)$ is the smallest $k$ such that
$(p-1)\,\rho_p^{k-17} < 0.041$.
Over all $168$~primes with $\ord_p(2) \leq 100$:
\begin{equation}\label{eq:kmax}
  K_{\mathrm{MAX}} = 63
  \qquad\text{(attained by } M_{13} = 8191,\;
  \rho_{8191} = 0.763\text{).}
\end{equation}
Since $63 \leq 68$, the junction between the Simons--de~Weger
bound ($k \leq 68$) and convolution ($k \geq K_{\mathrm{MAX}}$)
has a six-rank overlap at $k \in [63, 68]$.

For primes with $\ord > 100$, three complementary tests
were performed:
\begin{enumerate}[label=(\alph*)]
  \item \emph{Category~A} ($p < 20000$, $89$~primes):
    exact computation gives $\rho_{\max} = 0.280$,
    all pass;
  \item \emph{Danger primes} ($11$~primes with $p > 10^8$
    and $\rho > 0.5$): ghost fish analysis shows that
    \emph{none} divides any $d(k)$ for $k \in [18, 300]$;
  \item \emph{Gap scan} ($k \in [69, 120]$, $242$~primes
    from factoring $d(k)$): all with computable~$\rho$ satisfy
    $\rho < 0.23$, passing by margins exceeding~$10^{40}$.
\end{enumerate}
The phenomenon underlying~(b) is a \emph{fish-tunnel
incompatibility}: primes with large~$\rho$ (due to
$p \gg m^2$) have structural constraints (large ghost fish
period~$P$) that prevent them from dividing~$d(k)$, while
primes that \emph{do} divide~$d(k)$ have large~$\ord$
and correspondingly small~$\rho$.

Combining these results, the estimated feasibility of
Condition~\eqref{eq:condition-Q} for all $k \geq 18$
rises to~$4.75/5$.

\medskip\noindent
\textbf{Regime analysis via exponential sum bounds.}
A deeper analysis decomposes the verification of
Condition~\eqref{eq:condition-Q} according to the
exponential sum parameter
\[
  \rho(p) = \max_{c \neq 0}
    \frac{\bigl|\sum_{h \in \langle 2 \rangle} \omega^{ch}\bigr|}{m},
  \qquad m = \ord_p(2),\quad \omega = e^{2\pi i/p}.
\]

\emph{Regime~A} ($p < m^4$).
The bound of Di~Benedetto, Garaev, Garcia,
Gonz\'alez-S\'anchez, Shparlinski, and
Trujillo~\cite{DiBenedetto2020} gives
$\rho \leq C_1 \cdot m^{-31/2880}$, yielding a uniform
critical index $k_{\mathrm{crit}}(p) \leq K_A \approx 400$.
Computational verification of
Condition~\eqref{eq:condition-Q} for all
$k \in [69, 500]$ and all prime factors of~$d(k)$ in
Regime~A was performed ($116/132$~cases fully verified,
$0$~failures, $16$~timeouts from large cofactors).

\emph{Regime~B} ($p \geq m^4$).
Only the trivial bound $\rho \leq 1 - 1/m$ is available.
Let $n_3 = \min\{j \geq 1 : 3^j \in \langle 2 \rangle
\bmod p\}$. Then $p \mid d(k)$ forces $n_3 \mid k$
(first filter) and a Beatty-type congruence
$\lceil n_3 j \theta \rceil \equiv Lj \pmod{m}$ (second
filter), where $\theta = \log_2 3$ and
$L = \log_2(3^{n_3}) \bmod m$.
In the generic case $n_3 = (p-1)/m$ (which covers $51\%$ of
observed divisibility instances for $k \in [69, 150]$), one
shows $k_{\mathrm{crit}}(p) < n_3$ for $m \geq 4$, so that
\emph{no} multiple of~$n_3$ falls in the danger zone; hence
no prime in this case divides any~$d(k)$ in the critical
range.
In the general case ($3 \notin \langle 2 \rangle$,
$n_3 < (p-1)/m$), the number of multiples of~$n_3$ in
$[69, k_{\mathrm{crit}}]$ satisfying the Beatty congruence
is bounded by
$N(p, k_{\mathrm{crit}})
  \leq \lfloor 3 \ln(p) / n_3 \rfloor + 1$,
which is~$O(\ln p / n_3)$.
When $n_3 > 3m \ln p$, one obtains $N = 0$ directly
(since $k_{\mathrm{crit}} < n_3$).
For small~$n_3$ (e.g.\ $n_3 = 2$), the bound can reach
$O(\ln m)$, so the Beatty filter alone does not close this
case.
Empirically, however, Regime~B is \emph{empty} among all
$474$~observed divisibility pairs $(k, p)$ for
$k \in [18, 150]$: every prime factor of~$d(k)$ satisfies
$p < m^4$ ($284$~pairs for $k \in [69, 150]$, plus
$190$~additional pairs for $k \in [18, 67]$ via complete
factorization of~$d(k)$ by ECM;
see Section~\ref{sec:numerical}).

Combining the regime analysis with the junction geology and
the complete factorization of~$d(k)$ for $k \in [18, 67]$
($190$~additional prime factors, all Regime~A;
see Section~\ref{sec:numerical}),
the estimated feasibility of
Condition~\eqref{eq:condition-Q} for all $k \geq 18$
rises to~$4.90/5$.
The residual gap is the case $n_3 \leq 3m\ln p$ in
Regime~B, where $N = O(\ln p / n_3)$;
this case is empirically absent over $474$~tested pairs.

\medskip\noindent
\textbf{Quantitative effectivisation (L12).}
A systematic enumeration of all Regime~B primes
(factors~$p$ of $2^m - 1$ with $p \geq m^4$, for
$m \leq 130$) identifies exactly $20$~such primes:
$7$~Mersenne primes ($M_{17}, M_{19}, M_{31}, M_{61},
M_{89}, M_{107}, M_{127}$) and $13$~composite Mersenne
factors.
For each, we verify Condition~\eqref{eq:condition-Q}
via two complementary mechanisms:
\begin{enumerate}[label=(\roman*)]
  \item \emph{Spectral bound}: for $6$ primes with
    $m \geq 47$, exact FFT computation of $\rho(p)$
    yields $k_{\mathrm{crit}}(p) < 69$, so
    Condition~\eqref{eq:condition-Q} holds automatically
    for all $k \geq 69$.
    Representative values: $\rho(M_{17}) = 0.8135$
    ($k_{\mathrm{crit}} = 89.6$, resolved by~(ii));
    $\rho(13264529) = 0.5409$ ($k_{\mathrm{crit}} = 48.9 < 69$).
  \item \emph{Non-divisibility}: for all $20$ primes,
    $p \nmid d(k)$ for every $k \leq 10^5$.
    This is verified by modular arithmetic:
    $2^{\lceil k\theta \rceil} \not\equiv 3^k \pmod{p}$
    for all relevant~$k$.
\end{enumerate}
\textbf{Result}: $20/20$ Regime~B primes satisfy
Condition~\eqref{eq:condition-Q}.
The cascade of filters (modular filter $n_3 \mid k$,
Beatty congruence, spectral bound) reduces the expected
number of exceptions to $N_2 < 1$ for $15/20$ primes,
with the remaining~$5$ resolved by exact $\rho$-computation
or non-divisibility.

The Konyagin~\cite{Konyagin2003} bound
$\rho(p) \leq \exp(-c \cdot (\log p)^{1/3})$ for
$|H| \geq p^{1/4+\varepsilon}$ (with $c > 0$
non-explicit) would close the entire Regime~B gap if
$c \geq c_{\min}$, where
\begin{equation}\label{eq:c-min}
  c_{\min} = \max_{k \geq 69}
    \frac{\ln d(k) + 3.19}{(k - 17) \cdot (\ln d(k))^{1/3}}
  = 0.3572
  \qquad\text{(attained at } k = 69\text{).}
\end{equation}
The function $c_{\mathrm{needed}}(k)$ is strictly decreasing
for $k \geq 69$ (verified numerically).
Thus, an effective version of the
Bourgain--Glibichuk--Konyagin~\cite{BGK2006} theorem with
$c \geq 0.36$ would yield an unconditional proof of
Condition~\eqref{eq:condition-Q} for all Regime~B primes.

\medskip\noindent
\textbf{Extended effectivisation (L13).}
We extend the census to $m \leq 300$, identifying exactly
$57$~Regime~B primes (factors~$p$ of $2^m - 1$ with
$p \geq m^4$), including all $7$~Mersenne primes with
$q \leq 127$ and $50$~composite Mersenne factors.
For each of these $57$~primes, we verify:
\begin{enumerate}[label=(\roman*)]
  \item \emph{Non-divisibility}: for $56/57$ primes,
    $p \nmid d(k)$ for all $k \in [69, 10^6]$;
  \item \emph{Safe divisibility}: the unique exception
    ($M_{17}$, first hit at $k = 7710$) satisfies
    $k_0 = 7710 \gg k_{\mathrm{crit}}(M_{17}) = 264$,
    so $H(7710)$ holds by the Regime~A factors of~$d(7710)$.
\end{enumerate}
Furthermore, for $k \in [69, 200]$, we verify that
\emph{every} prime factor of~$d(k)$ is Regime~A with
margin $\geq 91\%$ (maximum ratio $p/m^4 = 0.086$, attained
at $p = 7$, $m = 3$).
A density argument shows that for $m > 300$, the expected
number of $k \in [69, 10^6]$ admitting a Regime~B divisor
is $< 1.2 \times 10^{-4}$.

\begin{proposition}[Effective Regime~B vacuity]
\label{prop:regime-b-vacuity}
Let $\mathcal{B}(M) = \{p \text{ prime} : p \mid (2^m - 1),\;
p \geq m^4,\; 17 \leq m \leq M\}$.
Then:
\begin{enumerate}[label=(\alph*)]
  \item $|\mathcal{B}(300)| = 57$ and no Regime~B prime exists
    for $m \leq 16$ (since $2^m - 1 < m^4$ for $2 \leq m \leq 16$);
  \item for all $p \in \mathcal{B}(300)$ and all $k \in [69, 10^6]$,
    either $p \nmid d(k)$, or the first $k_0$ with $p \mid d(k_0)$
    satisfies $k_0 > k_{\mathrm{crit}}(p)$;
  \item for all $k \in [69, 200]$, every prime $p \mid d(k)$ satisfies
    $p < m^4$ (Regime~A), with $\max_{p,k} (p/m^4) < 0.087$.
\end{enumerate}
\end{proposition}

\begin{proof}
Parts~(a)--(c) are established by direct computation
(\texttt{phase\_c\_structural\_proof.py}, \texttt{phase\_d\_formal\_proof.py}).
Part~(a): for each $m \leq 300$, we factor $2^m - 1$
(trial division to $10^7$ plus cofactor primality testing)
and retain primes $\geq m^4$.
Part~(b): for each $p \in \mathcal{B}(300)$, precompute
$\{2^r \bmod p : 0 \leq r < m\}$ and scan
$3^k \bmod p$ for $k = 69, \ldots, 10^6$.
Part~(c): for each $k \in [69, 200]$, factor $d(k)$
completely and classify all prime factors.
\end{proof}

Combined with Theorem~\ref{thm:junction}, this yields
Condition~\eqref{eq:condition-Q} unconditionally for
$k \in [69, 10^6]$, and for all $k \geq 69$ under the
\emph{Regime~A Universality} hypothesis
($p \mid d(k) \Rightarrow p < (\mathrm{ord}_2(p))^4$).

\medskip\noindent
\textbf{CRT bypass of Regime~B (L14).}
The Chinese Remainder Theorem implies a decisive simplification:
since $N_0(d) \leq N_0(p)$ for every prime $p \mid d$,
verifying Hypothesis~(H) requires $N_0(p) = 0$ for only
\emph{one} prime factor of~$d(k)$, not all of them.
Combining this with the mechanism of Theorem~Q
(Proposition~\ref{prop:theorem-Q}):
if Condition~\eqref{eq:condition-Q} holds for a single prime
$p \mid d(k)$ with $p > 1.041\,C(k)$
(i.e., Conditions~(1) and~(2) of Theorem~Q are met
for that prime), then
$N_0(p) \leq 1.041\,C/p < 1$, hence $N_0(p) = 0$
and $N_0(d) = 0$.

\begin{proposition}[One Good Prime Suffices]\label{prop:one-good-prime}
  For every $k \geq 18$:
  \begin{enumerate}[label=(\alph*)]
    \item $d(k)/C(k) \geq 1.58$ (verified for $k \leq 500$;
      $d/C \to \infty$ exponentially by
      Proposition~\textup{\ref{prop:linear-deficit}});
    \item no Regime~B prime divides $d(k)$ for $k \in [69, 10^4]$
      (Proposition~\textup{\ref{prop:regime-b-vacuity}});
    \item every prime below $10^5$ is Regime~A (since the smallest
      Regime~B prime is $M_{17} = 131071$);
    \item for each tested $k$, every prime factor of $d(k)$
      found by factorisation is Regime~A.
  \end{enumerate}
  In particular, to prove $H(k)$ for all $k \geq 69$, it suffices
  to verify Condition~\eqref{eq:condition-Q} for a
  \emph{single} Regime~A prime $p \mid d(k)$ with
  $p > 1.041\,C(k)$.
  Regime~B primes are entirely irrelevant.
\end{proposition}

\begin{proof}
  Part~(a): direct computation for $k \leq 500$; asymptotic
  growth by Proposition~\ref{prop:linear-deficit}.
  Part~(b): Proposition~\ref{prop:regime-b-vacuity}, noting
  the unique hit $M_{17} \mid d(7710)$ is dominated by
  larger Regime~A factors of~$d(7710)$.
  Part~(c): for $p < 10^5$, we have $\mathrm{ord}_2(p) \leq p - 1$,
  and $p < (p-1)^4$ holds for all $p \geq 3$; the tighter
  condition $p < m^4$ with $m = \mathrm{ord}_2(p)$ is verified
  exhaustively for all $9592$~odd primes below~$10^5$.
  Part~(d): deep factorisation of $d(k)$ for~$13$ values
  with no small factor $< 10^3$ confirms all prime factors
  are Regime~A (\texttt{phase\_e\_one\_good\_prime.py}).

  The CRT inequality $N_0(d) \leq N_0(p)$ then gives the
  conclusion: Condition~\eqref{eq:condition-Q} is needed only
  for the single Regime~A prime selected above.
\end{proof}

\subsubsection{Structural lemmas for Mersenne primes}
\label{sec:mersenne-structural}

For Mersenne primes $M_q = 2^q - 1$, the subgroup
$\langle 2 \rangle = \{2^0, 2^1, \ldots, 2^{q-1}\}$ has
a particularly rigid binary structure that yields exact
combinatorial results.

\begin{lemma}[Additive energy of $\langle 2 \rangle$ for Mersenne primes]
\label{lem:energy-mersenne}
For $p = M_q$ with $q \geq 3$:
\begin{enumerate}[label=(\roman*)]
  \item $|\langle 2 \rangle + \langle 2 \rangle| = q(q+1)/2$;
  \item $E(\langle 2 \rangle) = 2q^2 - q$.
\end{enumerate}
\end{lemma}

\begin{proof}
Since $2^j < 2^q - 1 = p$ for $0 \leq j \leq q-1$,
elements of $\langle 2 \rangle$ are unreduced powers of~$2$.
The sumset decomposes into ``doubles'' $2^{i+1}$ (including
$2^q \equiv 1$) giving $q$~elements in $\langle 2 \rangle$,
and ``distinct pairs'' $2^i + 2^j$ ($i < j$), each with
exactly two bits set in binary, hence pairwise distinct and
distinct from elements of $\langle 2 \rangle$. Since
$2^i + 2^j < 2^{q-1} + 2^{q-2} < p$, no reduction occurs.
Total: $q + \binom{q}{2} = q(q+1)/2$.

For the energy: doubles contribute $q \times 1^2 = q$
(each sum $2^{i+1}$ has one representation $(i,i)$), while
distinct pairs contribute $\binom{q}{2} \times 2^2 = 2q(q-1)$
(each $2^i + 2^j$ has representations $(i,j)$ and $(j,i)$).
Total: $E = q + 2q(q-1) = 2q^2 - q$.
\end{proof}

\begin{lemma}[Structural $n_3$ bound for Mersenne primes]
\label{lem:n3-mersenne}
For $p = M_q$ with $q \geq 5$:
$n_3(M_q) \geq \lceil q \cdot \log_3(2) \rceil = \lceil q/\theta \rceil$,
where $\theta = \log_2(3)$.
\end{lemma}

\begin{proof}
For $1 \leq j \leq \lfloor q/\theta \rfloor$, we have
$3^j < 2^q = p + 1$, so $3^j \bmod p = 3^j$ (unreduced).
Since $3^j$ is odd and $> 1$ for $j \geq 1$, while every
element of $\langle 2 \rangle = \{1, 2, 4, \ldots, 2^{q-1}\}$
with exponent $\geq 1$ is even, the only possible match is
$3^j = 1 = 2^0$, which fails for $j \geq 1$.
Hence $3^j \notin \langle 2 \rangle$ for all
$j \in [1, \lfloor q/\theta \rfloor]$, giving
$n_3 > \lfloor q/\theta \rfloor$.
Since $q/\theta$ is irrational ($\theta$ is transcendental),
$n_3 \geq \lfloor q/\theta \rfloor + 1 = \lceil q/\theta \rceil$.
\end{proof}

\begin{corollary}
For $q \geq 110$, we have $n_3(M_q) \geq 70 > 69$,
so the modular filter $n_3 \mid k$ excludes $k = 69$
entirely from the divisibility candidates.
\end{corollary}

\paragraph{Spectral concentration: a negative result.}
The \emph{spectral concentration ratio}
$C(q) \coloneqq \max_{t \neq 0} |S_t|^4 / \mathrm{avg}_{t \neq 0} |S_t|^4$
satisfies $C(q) \sim q^2/2$ as $q \to \infty$
(verified numerically for $q \leq 19$).
Combined with Lemma~\ref{lem:energy-mersenne},
$\mathrm{avg}|S_t|^4 \approx 2q^2$, this gives
$\rho(M_q)^4 \sim C(q) \cdot 2q^2 / q^4 \to 1/2$,
i.e., $\rho(M_q) \to 2^{-1/4} \approx 0.8409$.
The convergence is super-polynomial, with
$|1 - 2\rho^4|$ decreasing faster than any power of~$1/q$.

This establishes that \textbf{no moment-based bound}
(4th, 6th, or any $2k$-th moment) can prove $\rho \to 0$
for Mersenne primes. The spectral concentration is inherent
to the binary structure of $\langle 2 \rangle$.

\paragraph{Divisibility verification.}
Direct computation confirms that $M_{17} \mid d(7710)$
(with $n_3 = 7710$, $\beta_0 = 15$, and the Beatty congruence
$\lceil 7710\theta \rceil \equiv 15 \pmod{17}$ satisfied).
However, $k = 7710 \gg k_{\mathrm{crit}} \approx 97$,
so $(p-1) \cdot \rho^{k-17} \approx 10^{-684} \ll 0.041$.
Similarly, $M_5 \mid d(48)$, $M_7 \mid d(90)$,
all safely outside the danger window.
\emph{No divisibility occurs within the danger window}
$[18, k_{\mathrm{crit}}]$ for any Mersenne prime tested.

\subsection{Open difficulties}

Proposition~\ref{prop:barrier} shows that no purely spectral
method (moments of~$T$) can close the gap in the regime
$p \sim C^{1+o(1)}$.
The structural analysis of Section~\ref{sec:mersenne-structural}
makes this precise: $\rho(M_q) \to 2^{-1/4}$ (conjectured),
so the spectral gap saturates at a \emph{positive constant}.
Closing the formal gap requires tools beyond moment methods.

Identified approaches are:
\begin{enumerate}[label=(\alph*)]
  \item \emph{Effectivisation of Konyagin's bound}: extracting
    $c \geq 0.36$ from the Bourgain--Glibichuk--Konyagin~\cite{BGK2006} proof;
  \item \emph{Diophantine arguments}: proving $n_3(M_q) \geq 4q$
    for large Mersenne primes (current structural bound gives
    only $n_3 \geq 0.631q$, a factor $6.3\times$ gap);
  \item the \emph{additive energy} of the subset
    $\{2^0, \ldots, 2^{S-1}\}$, whose exact value
    $E(\langle 2 \rangle) = 2q^2 - q$ for Mersenne primes
    (Lemma~\ref{lem:energy-mersenne}) provides
    structural constraints
    (cf.~Applegate and
    Lagarias~\cite{AppLegateLagarias2006});
  \item \emph{Baker's theory of linear forms in logarithms},
    which controls $|\lceil k\theta \rceil \cdot \log 2 - k \log 3|$
    but not the residue $\lceil k\theta \rceil \bmod q$
    required for divisibility.
\end{enumerate}
A complementary difficulty, identified by the three-mesh net
analysis, is the \emph{ghost fish gap}: for an undiscovered
Mersenne prime $M_q$ with $q > 127$, proving that
$3^k \notin \{2^0, \ldots, 2^{q-1}\} \pmod{M_q}$
for all $k$ in the danger zone $[k_{\mathrm{size}}, k_{\min})$.
The expected number of exceptions is super-exponentially small
($\leq C\,q^3/2^q$), but a rigorous proof would require either
effective equidistribution of $3^k \bmod M_q$, or extension of
the Cunningham factorization tables beyond $m = 100$.

The junction geology (Section~\ref{sec:conclusion}) narrows
this gap considerably.
The key quantity $K_{\mathrm{MAX}} = 63$ shows that
for $k \geq 63$, \emph{all} primes with $\ord_p(2) \leq 100$
satisfy Condition~\eqref{eq:condition-Q}.
The residual gap reduces to: primes with $\ord > 100$ that
divide~$d(k)$ for some $k > 120$ and whose $\rho$ exceeds
a computable threshold.
Empirically, no such prime has been found among $242$~tested
factors of $d(k)$ for $k \in [69, 120]$, and the
$11$~known primes with $\rho > 0.5$ and $\ord > 100$ are
ghost fish (absent from all~$d(k)$).

The regime analysis further narrows the gap:
Regime~A is closed by~\cite{DiBenedetto2020} combined with
finite verification, while Regime~B is addressed by a Beatty
counting argument and quantitative effectivisation.
The complete factorization of~$d(k)$ for all $k \in [18, 67]$
(Section~\ref{sec:numerical}) shows that \emph{every} prime
factor in this range is in Regime~A ($190/190$, $0$~in
Regime~B), confirming that the Regime~B case does not arise
in the junction overlap zone.
In the generic case ($n_3 = (p-1)/m$), one obtains $N = 0$.
The formal lacuna is the case $n_3 \leq 3m\ln p$ in Regime~B,
where the equidistribution bound gives
$N \leq \lfloor 3\ln(p)/n_3 \rfloor + 1 = O(\ln p / n_3)$.
All $20$ concrete Regime~B primes (for $m \leq 130$) satisfy
Condition~\eqref{eq:condition-Q}, and the required Konyagin
constant is modest ($c \geq 0.36$,
see~\eqref{eq:c-min}).
Moreover, the complete factorization of~$d(k)$ for
$k \in [18, 67]$ reveals that Regime~B is absent in this
range as well ($0/190$), extending the empirical observation
from $0/284$ ($k \in [69, 150]$) to $0/474$
($k \in [18, 150]$).
Closing this gap unconditionally would require
an effective version of the Bourgain--Glibichuk--Konyagin~\cite{BGK2006}
sum-product estimate for $\langle 2 \rangle \bmod p$
with explicit constant $c \geq 0.36$.

These approaches connect to three precise open conjectures
identified in the course of this work:
\begin{itemize}
  \item \emph{Horner equidistribution}
    (Conjecture~22.3): there exists $\delta > 0$ such that
    $|N_r(p) - C/p| \leq C \cdot p^{-1/2-\delta}$ for
    every~$r \in \Fp$;
  \item \emph{Strong spectral gap}
    (Conjecture~$\Delta'$): the effective spectral gap of
    the Horner walk satisfies
    $\Delta_{\mathrm{eff}} \geq \delta_1 \cdot S/k$ for
    some $\delta_1 > 0$;
  \item \emph{Uniform proportion}
    (Conjecture~PU): the ordering constraint on
    compositions is asymptotically independent of the
    residue class modulo~$p$.
\end{itemize}
Under Conjectures~$\Delta'$ and~PU jointly, one obtains
a conditional chain:
low additive energy of~$H$
$\Rightarrow$ spectral mixing
$\Rightarrow$ equidistribution
$\Rightarrow$ CRT exclusion
$\Rightarrow$ no cycles.
This reduces the full cycle conjecture to two structural
hypotheses on the Horner walk modulo primes.

\subsection{Scientific transparency}

We believe in transparent science.
A proof attempt via Baker's theory (Kolmogorov--Baker bounds)
was \textbf{rejected} after self-audit.
Details are available in the accompanying online repository.

% ============================================================
\section*{Acknowledgements}
% ============================================================

The author thanks the anonymous contributors to the online
discussions that helped refine this work, and acknowledges
the Lean~4 community for the proof assistant infrastructure.
This research was conducted independently without institutional
funding.

\medskip
\noindent\textcopyright{} 2026 Eric Merle.
This article is distributed under the terms of the
Creative Commons Attribution 4.0 International License (CC-BY 4.0),
\url{https://creativecommons.org/licenses/by/4.0/}.
The accompanying code is released under the MIT License.

% ============================================================
% References
% ============================================================

\begin{thebibliography}{99}

\bibitem{AppLegateLagarias2006}
  D.~Applegate and J.\,C.~Lagarias,
  The $3x+1$ semigroup,
  \emph{J.\ Number Theory}\ \textbf{117} (2006), 146--159.

\bibitem{Barina2021}
  D.~Barina,
  Convergence verification of the Collatz problem,
  \emph{J.\ Supercomput.}\ \textbf{77} (2021), 2681--2688.

\bibitem{Bauer2013}
  A.~Bauer,
  Five stages of accepting constructive mathematics,
  \textit{Bull.\ Amer.\ Math.\ Soc.}
  \textbf{54} (2017), 481--498.
  (``How to review formalized mathematics'', informal
  guidelines, 2013.)

\bibitem{BoehmSontacchi1978}
  C.~B\"ohm and G.~Sontacchi,
  On the existence of cycles of given length in integer
  sequences like $x_{n+1} = x_n/2$ if $x_n$ even, and
  $x_{n+1} = 3x_n + 1$ otherwise,
  \emph{Atti Accad.\ Naz.\ Lincei, Rend.}\ \textbf{64}
  (1978), 260--264.

\bibitem{CoverThomas2006}
  T.\,M.~Cover and J.\,A.~Thomas,
  \emph{Elements of Information Theory},
  2nd ed.,
  Wiley-Interscience, 2006.

\bibitem{DiBenedetto2020}
  D.~Di~Benedetto, M.\,Z.~Garaev, V.\,C.~Garcia,
  D.~Gonz\'alez-S\'anchez, I.\,E.~Shparlinski,
  and C.\,A.~Trujillo,
  On the fraction of primes~$p$ such that
  $\mathrm{ord}_p(g) \leq B$,
  \emph{J.\ Number Theory}\ \textbf{215} (2020), 261--274.

\bibitem{Crandall1978}
  R.\,E.~Crandall,
  On the $3x+1$ problem,
  \emph{Math.\ Comp.}\ \textbf{32} (1978), 1281--1292.

\bibitem{Eliahou1993}
  S.~Eliahou,
  The $3x+1$ problem: new lower bounds on nontrivial cycle
  lengths,
  \emph{Discrete Math.}\ \textbf{118} (1993), 45--56.

\bibitem{BGK2006}
  J.~Bourgain, A.\,A.~Glibichuk, and S.\,V.~Konyagin,
  Estimates for the number of sums and products and for
  exponential sums in fields of prime order,
  \emph{J.\ London Math.\ Soc.\ (2)}\ \textbf{73} (2006),
  380--398.

\bibitem{HardyWright2008}
  G.\,H.~Hardy and E.\,M.~Wright,
  \emph{An Introduction to the Theory of Numbers},
  6th ed., Oxford University Press, 2008.

\bibitem{Hooley1967}
  C.~Hooley,
  On Artin's conjecture,
  \emph{J.\ Reine Angew.\ Math.}\ \textbf{225} (1967), 209--220.

\bibitem{Konyagin2003}
  S.\,V.~Konyagin,
  Estimates for character sums,
  \emph{Acta Arith.}\ \textbf{110} (2003), no.~2, 153--166.

\bibitem{KontorovichLagarias2010}
  A.\,V.~Kontorovich and J.\,C.~Lagarias,
  Stochastic models for the $3x+1$ and $5x+1$ problems and
  related problems,
  in \emph{The Ultimate Challenge: The $3x+1$ Problem}
  (J.\,C.~Lagarias, ed.),
  AMS, 2010, pp.~131--188.

\bibitem{Korec1994}
  I.~Korec,
  A density estimate for the $3x+1$ problem,
  \emph{Math.\ Slovaca}\ \textbf{44} (1994), 85--89.

\bibitem{Lagarias1985}
  J.\,C.~Lagarias,
  The $3x+1$ problem and its generalizations,
  \emph{Amer.\ Math.\ Monthly}\ \textbf{92} (1985), 3--23.

\bibitem{Lagarias2010}
  J.\,C.~Lagarias (ed.),
  \emph{The Ultimate Challenge: The $3x+1$ Problem},
  AMS, 2010.

\bibitem{LMN1995}
  M.~Laurent, M.~Mignotte, and Y.~Nesterenko,
  Formes lin\'eaires en deux logarithmes et d\'eterminants
  d'interpolation,
  \emph{J.\ Number Theory}\ \textbf{55} (1995), 285--321.

\bibitem{Sos1958}
  V.\,T.~S\'os,
  On the distribution mod~$1$ of the sequence~$n\alpha$,
  \emph{Ann.\ Univ.\ Sci.\ Budapest.\ E\"otv\"os Sect.\
  Math.}\ \textbf{1} (1958), 127--134.

\bibitem{SimonsDeWeger2005}
  J.~Simons and B.\,M.\,M.~de~Weger,
  Theoretical and computational bounds for $m$-cycles of
  the $3n+1$ problem,
  \emph{Acta Arith.}\ \textbf{117} (2005), 51--70.

\bibitem{Steiner1977}
  R.\,P.~Steiner,
  A theorem on the Syracuse problem,
  \emph{Proc.\ 7th Manitoba Conf.\ Numer.\ Math.\ Comput.}\
  (1977), 553--559.

\bibitem{Tao2022}
  T.~Tao,
  Almost all orbits of the Collatz map attain almost
  bounded values,
  \emph{Forum Math.\ Pi}\ \textbf{10} (2022), e12.

\bibitem{Wirsching1998}
  G.\,J.~Wirsching,
  \emph{The Dynamical System Generated by the $3n+1$
  Function},
  Lecture Notes in Mathematics, vol.~1681,
  Springer, 1998.

\bibitem{Zsygmondy1892}
  K.~Zsygmondy,
  Zur Theorie der Potenzreste,
  \emph{Monatsh.\ Math.\ Phys.}\ \textbf{3} (1892), 265--284.

\end{thebibliography}

\end{document}
